\documentclass[11pt]{amsart}
\usepackage[margin=1in]{geometry}
\usepackage[T1]{fontenc}
\usepackage{amsmath,amssymb,amsthm,mathrsfs}
\usepackage{booktabs,microtype}
\usepackage{tcolorbox}
\usepackage{hyperref}
\hypersetup{hidelinks}
\allowdisplaybreaks[2]
\raggedbottom
\setlength{\emergencystretch}{2em}

\newtheorem{theorem}{Theorem}[section]
\newtheorem{lemma}[theorem]{Lemma}
\newtheorem{proposition}[theorem]{Proposition}
\newtheorem{corollary}[theorem]{Corollary}
\theoremstyle{remark}
\newtheorem{remark}[theorem]{Remark}
\theoremstyle{plain}
\newtheorem*{valueagain}{Theorem \ref{thm:value} (restated)}
\newtheorem*{robustagain}{Theorem \ref{thm:robust} (restated)}
\newtheorem*{noiseagain}{Theorem \ref{thm:noise} (restated)}
\newtheorem*{supremumagain}{Theorem \ref{thm:supremum} (restated)}
\numberwithin{equation}{section}

\DeclareMathOperator{\Tr}{Tr}
\DeclareMathOperator{\rank}{rank}
\DeclareMathOperator{\diag}{diag}
\DeclareMathOperator{\spec}{spec}
\DeclareMathOperator{\LF}{LF}
\DeclareMathOperator{\TV}{TV}
\DeclareMathOperator{\Err}{Err}
\DeclareMathOperator{\Beta}{Beta}
\DeclareMathOperator{\Bin}{Bin}
\DeclareMathOperator{\sech}{sech}
\DeclareMathOperator{\artanh}{artanh}
\DeclareMathOperator{\amag}{amag}
\DeclareMathOperator{\BC}{BC}
\newcommand{\C}{\mathbb C}
\newcommand{\R}{\mathbb R}
\newcommand{\Z}{\mathbb Z}
\newcommand{\ii}{\mathrm i}
\newcommand{\dd}{\,\mathrm d}
\newcommand{\ind}{\mathbf 1}
\newcommand{\Loc}{\mathcal C}
\newcommand{\NS}{\mathcal P}
\newcommand{\ket}[1]{\lvert#1\rangle}
\newcommand{\bra}[1]{\langle#1\rvert}
\newcommand{\abs}[1]{\left|#1\right|}
\newcommand{\norm}[1]{\left\|#1\right\|}
\newcommand{\ip}[2]{\langle#1,#2\rangle}
\newcommand{\Gc}{G_{\rm Cat}}
\newcommand{\Fdef}{\mathfrak F_d}
\newcommand{\Gdef}{\mathcal G_d}
\newcommand{\Ddef}{\mathcal D_d}
\newcommand{\IF}{I_d^{\rm F}}
\newcommand{\cstar}{c_*}
\newcommand{\cKL}{c_{\rm KL}}
\newcommand{\Hilb}{\mathsf H}
\newcommand{\Comp}[1]{\mathsf C^{(#1)}}
\newcommand{\RootP}{\mathcal M}
\newcommand{\epsanc}{\varepsilon_{\rm anc}}
\newcommand{\epsred}{\varepsilon_{\rm red}}
\newcommand{\epsmeas}{\varepsilon_{\rm meas}}
\newcommand{\epssc}{\varepsilon_{\rm sc}}
\newcommand{\doi}[1]{\href{https://doi.org/#1}{\nolinkurl{doi:#1}}}
\newcommand{\arxiv}[1]{\href{https://arxiv.org/abs/#1}{\nolinkurl{arXiv:#1}}}

\begin{filecontents*}[overwrite]{cglmp-references.bib}
@article{CGLMP,
 author = {D. Collins and N. Gisin and N. Linden and S. Massar and S. Popescu},
 title = {Bell inequalities for arbitrarily high-dimensional systems},
 journal = {Phys. Rev. Lett.},
 volume = {88},
 pages = {040404},
 year = {2002},
 doi = {10.1103/PhysRevLett.88.040404},
 eprint = {quant-ph/0106024},
 note = {\doi{10.1103/PhysRevLett.88.040404}; \arxiv{quant-ph/0106024}}
}
@article{DKZ,
 author = {T. Durt and D. Kaszlikowski and M. {\.Z}ukowski},
 title = {Violations of local realism with quantum systems described by
          {$N$}-dimensional {Hilbert} spaces up to {$N=16$}},
 journal = {Phys. Rev. A},
 volume = {64},
 pages = {024101},
 year = {2001},
 doi = {10.1103/PhysRevA.64.024101},
 eprint = {quant-ph/0101084},
 note = {\doi{10.1103/PhysRevA.64.024101}; \arxiv{quant-ph/0101084}}
}
@article{ADGL,
 author = {A. Ac{\'i}n and T. Durt and N. Gisin and J. I. Latorre},
 title = {Quantum nonlocality in two three-level systems},
 journal = {Phys. Rev. A},
 volume = {65},
 pages = {052325},
 year = {2002},
 doi = {10.1103/PhysRevA.65.052325},
 eprint = {quant-ph/0111143},
 note = {\doi{10.1103/PhysRevA.65.052325}; \arxiv{quant-ph/0111143}}
}
@article{Masanes,
 author = {Ll. Masanes},
 title = {Tight {Bell} inequality for {$d$}-outcome measurements correlations},
 journal = {Quantum Inf. Comput.},
 volume = {3},
 number = {4},
 pages = {345--358},
 year = {2003},
 doi = {10.26421/QIC3.4-4},
 eprint = {quant-ph/0210073},
 note = {\doi{10.26421/QIC3.4-4}; \arxiv{quant-ph/0210073}}
}
@article{VDGG,
 author = {W. van Dam and R. D. Gill and P. D. Gr{\"u}nwald},
 title = {The statistical strength of nonlocality proofs},
 journal = {IEEE Trans. Inform. Theory},
 volume = {51},
 number = {8},
 pages = {2812--2835},
 year = {2005},
 doi = {10.1109/TIT.2005.851738},
 eprint = {quant-ph/0307125},
 note = {\doi{10.1109/TIT.2005.851738}; \arxiv{quant-ph/0307125}}
}
@article{AGG,
 author = {A. Ac{\'i}n and R. Gill and N. Gisin},
 title = {Optimal {Bell} tests do not require maximally entangled states},
 journal = {Phys. Rev. Lett.},
 volume = {95},
 pages = {210402},
 year = {2005},
 doi = {10.1103/PhysRevLett.95.210402},
 eprint = {quant-ph/0506225},
 note = {\doi{10.1103/PhysRevLett.95.210402}; \arxiv{quant-ph/0506225}}
}
@article{Fine,
 author = {A. Fine},
 title = {Hidden variables, joint probability, and the {Bell} inequalities},
 journal = {Phys. Rev. Lett.},
 volume = {48},
 pages = {291--295},
 year = {1982},
 doi = {10.1103/PhysRevLett.48.291},
 note = {\doi{10.1103/PhysRevLett.48.291}}
}
@article{CG,
 author = {D. Collins and N. Gisin},
 title = {A relevant two qubit {Bell} inequality inequivalent to the {CHSH} inequality},
 journal = {J. Phys. A: Math. Gen.},
 volume = {37},
 pages = {1775--1787},
 year = {2004},
 doi = {10.1088/0305-4470/37/5/021},
 eprint = {quant-ph/0306129},
 note = {\doi{10.1088/0305-4470/37/5/021}; \arxiv{quant-ph/0306129}}
}
@article{BGP,
 author = {J.-D. Bancal and N. Gisin and S. Pironio},
 title = {Looking for symmetric {Bell} inequalities},
 journal = {J. Phys. A: Math. Theor.},
 volume = {43},
 pages = {385303},
 year = {2010},
 doi = {10.1088/1751-8113/43/38/385303},
 eprint = {1004.4146},
 note = {\doi{10.1088/1751-8113/43/38/385303}; \arxiv{1004.4146}}
}
@article{JC,
 author = {J. Jesus and E. {Zambrini Cruzeiro}},
 title = {Tight {Bell} inequalities from polytope slices},
 journal = {Phys. Rev. A},
 volume = {108},
 pages = {052220},
 year = {2023},
 doi = {10.1103/PhysRevA.108.052220},
 eprint = {2212.03212},
 note = {\doi{10.1103/PhysRevA.108.052220}; \arxiv{2212.03212}}
}
@article{LVN,
 author = {B. Lang and T. V{\'e}rtesi and M. Navascu{\'e}s},
 title = {Closed sets of correlations: answers from the zoo},
 journal = {J. Phys. A: Math. Theor.},
 volume = {47},
 pages = {424029},
 year = {2014},
 doi = {10.1088/1751-8113/47/42/424029},
 eprint = {1402.2850},
 note = {\doi{10.1088/1751-8113/47/42/424029}; \arxiv{1402.2850}}
}
@article{NPA,
 author = {M. Navascu{\'e}s and S. Pironio and A. Ac{\'i}n},
 title = {Bounding the set of quantum correlations},
 journal = {Phys. Rev. Lett.},
 volume = {98},
 pages = {010401},
 year = {2007},
 doi = {10.1103/PhysRevLett.98.010401},
 eprint = {quant-ph/0607119},
 note = {\doi{10.1103/PhysRevLett.98.010401}; \arxiv{quant-ph/0607119}}
}
@misc{IR,
 author = {M. Ioannou and D. Rosset},
 title = {Noncommutative polynomial optimization under symmetry},
 howpublished = {arXiv preprint},
 year = {2021},
 eprint = {2112.10803},
 note = {\arxiv{2112.10803}}
}
@article{SATWAP,
 author = {A. Salavrakos and R. Augusiak and J. Tura and P. Wittek
           and A. Ac{\'i}n and S. Pironio},
 title = {Bell inequalities tailored to maximally entangled states},
 journal = {Phys. Rev. Lett.},
 volume = {119},
 pages = {040402},
 year = {2017},
 doi = {10.1103/PhysRevLett.119.040402},
 eprint = {1607.04578},
 note = {\doi{10.1103/PhysRevLett.119.040402}; \arxiv{1607.04578}}
}
@article{Kaniewski,
 author = {J. Kaniewski and I. {\v{S}}upi{\'c} and J. Tura and F. Baccari
           and A. Salavrakos and R. Augusiak},
 title = {Maximal nonlocality from maximal entanglement and mutually unbiased
          bases, and self-testing of two-qutrit quantum systems},
 journal = {Quantum},
 volume = {3},
 pages = {198},
 year = {2019},
 doi = {10.22331/q-2019-10-24-198},
 eprint = {1807.03332},
 note = {\doi{10.22331/q-2019-10-24-198}; \arxiv{1807.03332}}
}
@article{Farkas,
 author = {M. Farkas and J. Kaniewski},
 title = {Self-testing mutually unbiased bases in the prepare-and-measure scenario},
 journal = {Phys. Rev. A},
 volume = {99},
 pages = {032316},
 year = {2019},
 doi = {10.1103/PhysRevA.99.032316},
 eprint = {1803.00363},
 note = {\doi{10.1103/PhysRevA.99.032316}; \arxiv{1803.00363}}
}
@article{Yang,
 author = {T. H. Yang and T. V{\'e}rtesi and J.-D. Bancal and V. Scarani
           and M. Navascu{\'e}s},
 title = {Robust and versatile black-box certification of quantum devices},
 journal = {Phys. Rev. Lett.},
 volume = {113},
 pages = {040401},
 year = {2014},
 doi = {10.1103/PhysRevLett.113.040401},
 eprint = {1406.7127},
 note = {\doi{10.1103/PhysRevLett.113.040401}; \arxiv{1406.7127}}
}
@article{Sarkar,
 author = {S. Sarkar and D. Saha and J. Kaniewski and R. Augusiak},
 title = {Self-testing quantum systems of arbitrary local dimension with
          minimal number of measurements},
 journal = {npj Quantum Inf.},
 volume = {7},
 pages = {151},
 year = {2021},
 doi = {10.1038/s41534-021-00490-3},
 eprint = {1909.12722},
 note = {\doi{10.1038/s41534-021-00490-3}; \arxiv{1909.12722}}
}
@misc{Heinavaara,
 author = {O. Hein{\"a}vaara},
 title = {Convolution comparison measures},
 howpublished = {arXiv preprint},
 year = {2026},
 eprint = {2602.10373},
 note = {\arxiv{2602.10373}}
}
@article{ZG,
 author = {S. Zohren and R. D. Gill},
 title = {Maximal violation of the {CGLMP} inequality for infinite dimensional states},
 journal = {Phys. Rev. Lett.},
 volume = {100},
 pages = {120406},
 year = {2008},
 doi = {10.1103/PhysRevLett.100.120406},
 eprint = {quant-ph/0612020},
 note = {\doi{10.1103/PhysRevLett.100.120406}; \arxiv{quant-ph/0612020}}
}
@article{SCB,
 author = {I. {\v{S}}upi{\'c} and D. Cavalcanti and J. Bowles},
 title = {Device-independent certification of tensor products of quantum states
          using single-copy self-testing protocols},
 journal = {Quantum},
 volume = {5},
 pages = {418},
 year = {2021},
 doi = {10.22331/q-2021-03-23-418},
 eprint = {1909.12759},
 note = {\doi{10.22331/q-2021-03-23-418}; \arxiv{1909.12759}}
}
@article{Grudka,
 author = {A. Grudka and K. Horodecki and M. Horodecki and P. Horodecki
           and R. Horodecki and P. Joshi and W. K{\l}obus and A. W{\'o}jcik},
 title = {Quantifying contextuality},
 journal = {Phys. Rev. Lett.},
 volume = {112},
 pages = {120401},
 year = {2014},
 doi = {10.1103/PhysRevLett.112.120401},
 eprint = {1209.3745},
 note = {\doi{10.1103/PhysRevLett.112.120401}; \arxiv{1209.3745}}
}
@article{Joshi,
 author = {P. Joshi and K. Horodecki and M. Horodecki and P. Horodecki
           and R. Horodecki and B. Li and S. J. Szarek and T. Szarek},
 title = {Bound on {Bell} inequalities by fraction of determinism and reverse
          triangle inequality},
 journal = {Phys. Rev. A},
 volume = {92},
 pages = {032329},
 year = {2015},
 doi = {10.1103/PhysRevA.92.032329},
 eprint = {1502.03088},
 note = {\doi{10.1103/PhysRevA.92.032329}; \arxiv{1502.03088}}
}
@article{GMS,
 author = {N. Gisin and A. A. M{\'e}thot and V. Scarani},
 title = {Pseudo-telepathy: input cardinality and {Bell}-type inequalities},
 journal = {Int. J. Quantum Inf.},
 volume = {5},
 pages = {525--534},
 year = {2007},
 doi = {10.1142/S021974990700289X},
 eprint = {quant-ph/0610175},
 note = {\doi{10.1142/S021974990700289X}; \arxiv{quant-ph/0610175}}
}
@article{BRL,
 author = {K. Baek and J. Ryu and J. Lee},
 title = {Robustness measures for quantifying nonlocality},
 journal = {New J. Phys.},
 volume = {27},
 pages = {053001},
 year = {2025},
 doi = {10.1088/1367-2630/adcd5a},
 eprint = {2311.07077},
 note = {\doi{10.1088/1367-2630/adcd5a}; \arxiv{2311.07077}}
}
@misc{OQP,
 author = {{IQOQI Vienna}},
 title = {The power of {CGLMP} inequalities},
 howpublished = {Open Quantum Problems, Problem~27},
 year = {2026},
 note = {\url{https://oqp.iqoqi.oeaw.ac.at/the-power-of-cglmp-inequalities/}, accessed 8 September 2026}
}
@misc{Staufenbiel,
 author = {C. Staufenbiel},
 title = {Bell inequalities from polyhedral sampling},
 howpublished = {arXiv preprint},
 year = {2026},
 eprint = {2604.22859},
 note = {\arxiv{2604.22859}}
}
@inproceedings{FukudaProdon,
 author = {K. Fukuda and A. Prodon},
 title = {Double description method revisited},
 booktitle = {Combinatorics and Computer Science},
 editor = {M. Deza and R. Euler and I. Manoussakis},
 series = {Lecture Notes in Computer Science},
 volume = {1120},
 pages = {91--111},
 publisher = {Springer},
 address = {Berlin, Heidelberg},
 year = {1996},
 doi = {10.1007/3-540-61576-8_77},
 note = {\doi{10.1007/3-540-61576-8_77}}
}
\end{filecontents*}

\title[CGLMP bounds and statistical strength]
{CGLMP bounds, measurement rigidity, and statistical strength\\
in two-setting Bell tests}
\author{\ }
% Affiliation: \address{Department, institution, city, postal code, country}
% Correspondence: \email{address@example.org}
% Acknowledgements: \thanks{Acknowledgements to be supplied.}
% Funding: \thanks{Funding sources and grant numbers to be supplied.}
\date{}

\begin{document}

\begin{abstract}
We determine the sharp CGLMP bound for projective measurements on
maximally entangled states in the bipartite Bell scenario with two
settings and \(d\) outcomes; these questions form Problem~27 of the IQOQI list
of open problems in quantum information. Equality holds precisely when the local
dimension is a multiple of \(d\) and the projective measurements are
unitarily equivalent to Fourier-phase measurements with the
corresponding multiplicity. Quantitative estimates control the
measurement distance by the Bell-value deficit and give an arithmetic
dimension obstruction. The proof combines a matrix Hardy inequality
with a positive law on partitions and polynomial positivity bounds.
We then compare CGLMP optimization with noise robustness and
Kullback--Leibler statistical strength. Fourier-phase projective
measurements optimize uniform-noise tolerance for the CGLMP witness,
whereas two-setting tensor measurements achieve greater global
white-noise resistance and statistical strength. For every input
distribution, we determine the dimension-unrestricted
statistical-strength supremum and give a sequence approaching it.
A constructive local-component bound places every finite-outcome
quantum behaviour strictly below this supremum for full-support inputs.
\end{abstract}
\maketitle

\section{Introduction}\label{sec:introduction}

Bell experiments with two measurement choices and many outcomes offer
a direct way to study nonlocality in higher-dimensional quantum
systems. Increasing the number of outcomes makes the internal
structure of a qudit accessible while keeping the input alphabet
small. This combination is also useful for device-independent
certification, where a small number of measurement choices can
certify a large local Hilbert space \cite{SATWAP,Sarkar}. It motivates
the search for Bell inequalities whose optimal measurements admit
both an experimentally natural implementation and a mathematical
characterization.

The CGLMP inequalities of Collins, Gisin, Linden, Massar, and
Popescu \cite{CGLMP} provide a distinguished family in this setting.
Their local bound follows from a modular constraint around the four
measurement contexts, and each inequality defines a facet of the
local polytope \cite{Masanes}. The associated Fourier-phase
measurements combine a discrete Fourier transformation with diagonal
unitaries. Durt, Kaszlikowski, and \.Zukowski \cite{DKZ} conjectured
that these measurements are optimal on the maximally entangled
state. Zohren and Gill \cite{ZG} subsequently examined the same
measurement family in their analysis of large-dimensional CGLMP
violations. Beyond the optimal value, a natural rigidity question
asks whether every maximizing projective strategy has this form.

The fixed-state problem has a useful tracial structure: correlations
on a maximally entangled state are normalized traces of products of
measurement operators. Lang, V\'ertesi, and Navascu\'es \cite{LVN}
exploited this structure in semidefinite relaxations and obtained
sharp numerical evidence for the three-outcome value. Optimization
over states leads to a different problem. Ac\'in, Durt, Gisin, and
Latorre \cite{ADGL} showed that nonmaximally entangled states can give
larger CGLMP violations. The hierarchy of Navascu\'es, Pironio, and
Ac\'in \cite{NPA}, together with symmetry-adapted polynomial methods
\cite{IR}, provides arbitrary-state bounds, including tight analytic
certificates at \(d=3,4\). Numerical black-box certification of the
qutrit optimum was developed by Yang et al.\ \cite{Yang}. Related
rigidity results concern mutually unbiased bases in Bell scenarios,
with three settings in the qutrit case \cite{Kaniewski}, and in
prepare-and-measure scenarios \cite{Farkas}.

A second issue is the meaning of an optimal Bell experiment. A large
value of a chosen functional determines the amount of uniform noise
that this particular witness tolerates. Global white-noise resistance
instead depends on the intersection of the noise line with the whole
local polytope. Statistical strength is a third criterion: the
quantum-to-local Kullback--Leibler divergence quantifies the expected
evidence against local realism per complete trial \cite{VDGG,AGG}.
It depends on the entire behaviour and on the distribution of the
inputs. Ac\'in, Gill, and Gisin \cite{AGG} compared these choices in
the design of Bell tests, while Baek, Ryu, and Lee \cite{BRL}
studied distinctions between robustness measures. These considerations
make it important to keep the witness, the noise model, and the
statistical objective separate.

The questions studied in this paper were codified by Gill as Problem~27
of the IQOQI list of open problems in quantum information \cite{OQP}.
We reproduce its statement, with the citation markers of the original
resolved to the present bibliography.

\begin{tcolorbox}[colback=white,colframe=black!60,boxrule=0.5pt,arc=1.5pt,
  left=6pt,right=6pt,top=4pt,bottom=4pt,
  title={Problem 27 (IQOQI Open Quantum Problems): The power of CGLMP inequalities},
  coltitle=white,colbacktitle=black!60]
\small
In the setting of Problem~26, consider especially the case \((N,M,K)=(2,2,d)\).

\smallskip\noindent\textbf{Problem 27.\,A.}
Show that every face of the local polytope \(C\), which is not already
contained in a face of the no-signalling polytope \(P\) is of CGLMP type,
i.e., an inequality of the form first written out in \cite{CGLMP}, but
possibly lifted from lower dimensions by fusing together some outcomes.

\smallskip\noindent\textbf{Problem 27.\,B.}
Numerically, the observables maximally violating the CGLMP inequality on
a maximally entangled state are of a very specific form \cite{DKZ},
involving measurements in computational basis, transformed by only
discrete Fourier transformation and diagonal unitaries \cite{CGLMP}. Show
that this is necessarily the case. Show also that these measurements
realize the highest resistance of violation to noise, and the best
discrimination against classical realism in the sense of
Kullback--Leibler divergence \cite{VDGG}.
\end{tcolorbox}

Here \(N\), \(M\) and \(K\) denote the numbers of parties, settings per
party and outcomes per measurement, and Problem~26 of the same list
concerns Bell inequalities in this general setting. The answers are as
follows. Part~A fails, as Bancal, Gisin, and Pironio already observed
in answering Gill's question in the negative \cite{BGP}: their four-outcome
facet, recalled in Appendix~\ref{app:facet} (Proposition~\ref{prop:bgp})
with a self-contained certificate, lies in no face of the no-signalling
polytope and is neither a CGLMP inequality nor a lifting of one. In
Part~B, the first assertion holds for projective measurements:
Theorem~\ref{thm:value} shows that on a maximally entangled state the
CGLMP value is maximized only by Fourier-phase measurements, in every
outcome count, and Theorem~\ref{thm:robust} makes this quantitative.
The two optimality assertions of Part~B hold only in a restricted sense:
Theorem~\ref{thm:noise} shows that Fourier-phase measurements optimize
uniform-noise tolerance for the CGLMP witness itself but not global
white-noise resistance, and Theorem~\ref{thm:supremum} shows that they
do not optimize Kullback--Leibler statistical strength, whose
dimension-unrestricted supremum we determine for every input
distribution. At two outcomes, Fine's characterization gives the full
local description \cite{Fine}; higher-outcome facet enumeration already
exhibits a richer geometry \cite{CG,BGP,JC,Staufenbiel}.

\subsection*{Main results}

We normalize the CGLMP functional \(I_d\) to have local bound \(2\).
Write \(\Phi_n=n^{-1/2}\sum_{j=0}^{n-1}\ket{jj}\).
The standard Fourier-phase projectors \(\Pi^t_a\) are defined in
\eqref{eq:standard}. For the fixed-state statements, we use the tracial
ordering
\[
 (P^0,P^1,P^2,P^3)=(A_0^T,B_0,A_1^T,B_1).
\]

\begin{theorem}[Projective value and rigidity]\label{thm:value}
Let \(d\ge2\), \(n\ge1\), and let the four measurements on \(\Phi_n\)
be \(d\)-outcome PVMs. Then
\begin{equation}\label{eq:sharpvalue}
 I_d\le\IF:=
 \frac4{d(d-1)}\sum_{h=1}^{d-1}(d-h)
                  \sec\!\left(\frac{\pi h}{2d}\right).
\end{equation}
Equality holds if and only if \(n=dk\) and, after a common unitary
conjugation in the tracial representation, the four PVMs are
\(\Pi^t_a\otimes I_k\). In particular, all projective maximizers on
\(\Phi_d\) are rank-one Fourier-phase measurements.
\end{theorem}

Sarkar et al.\ \cite{Sarkar} obtained a device-independent self-test
of the same measurements for the inequalities of Salavrakos
et al.\ \cite{SATWAP}. Theorem~\ref{thm:value} is a rigidity statement
for projective measurements on the fixed maximally entangled state
and for the CGLMP functional itself.

The distance \(\Err_{d,n}\) and constants \(C_d^{\rm st},R_d^{\rm st},K_d\)
are defined in \eqref{eq:distance} and \eqref{eq:stabilityconstants}.

\begin{theorem}[Quantitative rigidity and arithmetic obstruction]
\label{thm:robust}
Let \(d\ge2\), \(n\ge1\), and let \(P\) be the tracial PVMs of
\(d\)-outcome projective measurements on \(\Phi_n\).
Put \(\Delta=\IF-I_d\). If \(d\mid n\), then
\begin{equation}\label{eq:robust}
 \Err_{d,n}(P)\le\min\{8,K_d\Delta\},
 \qquad K_d\le786432\,d^6.
\end{equation}
For arbitrary \(n\), with \(s=n\bmod d\),
\begin{equation}\label{eq:arithmeticgap}
 \Delta\ge
 \frac{s(d-s)}{3n\,d^4(d+1)C_d^{\rm st}}.
\end{equation}
If \(\mathcal I_{d,n}\) denotes the projective maximum in that
dimension, then for \(s>0\),
\begin{equation}\label{eq:gaporder}
 \frac{s(d-s)}{3n\,d^4(d+1)C_d^{\rm st}}
 \le\IF-\mathcal I_{d,n}
 \le\frac{s}{n}(\IF-2).
\end{equation}
\end{theorem}

The arithmetic obstruction distinguishes the ambient dimension from
the number of outcomes: equality requires \(d\mid n\), while
\eqref{eq:gaporder} controls the optimal deficit when \(d\) does not
divide \(n\).

Write \(Q_d\) for the Fourier-phase behaviour in \eqref{eq:Qd} and
\(T_m\) for the common-input tensor CHSH behaviour in
\eqref{eq:tensorbehaviour}. The quantity \(v_{\rm crit}(q)\) denotes
the critical visibility of \(q\) under uniform noise
(Section~\ref{sec:preliminaries}).

\begin{theorem}[Witness tolerance and global noise resistance]
\label{thm:noise}
For \(d\)-outcome projective measurements on \(\Phi_n\) producing
\(q\) with \(I_d(q)>2\),
\begin{equation}\label{eq:witnessnoise}
 v_{\rm CGLMP}(q)=\frac2{I_d(q)}\ge\frac2{\IF}.
\end{equation}
Equality holds precisely for the standard amplifications in
Theorem~\ref{thm:value}; in particular, the minimum is attained
whenever \(d\mid n\).
For every integer \(m\ge20000\), the two rank-one measurements per
party defining \(T_m\) on \(\Phi_{2^m}\) satisfy
\begin{equation}\label{eq:noiseseparation}
 v_{\rm crit}(T_m)\le\frac{17}{29}
       <\frac35<v_{\rm crit}(Q_{2^m}).
\end{equation}
\end{theorem}

The comparison keeps the state, local dimension, number of settings,
and measurement ranks fixed. Common-input tensor measurements also
appear in the certification of tensor products with an unchanged
number of measurement choices \cite{SCB}.

For an input law \(\mu\), let \(\mathcal S_\mu\) be the statistical
strength in \eqref{eq:strength}, and let \(\Xi(\mu)\) be the
four-variable minimum in \eqref{eq:Xi}.

\begin{theorem}[Dimension-unrestricted statistical strength]
\label{thm:supremum}
For every fixed input distribution independent of the hidden variable,
\begin{equation}\label{eq:supremum}
 \sup_{\substack{d\ge2\\
     \text{two rank-one PVMs per party on }\Phi_d}}
       \mathcal S_\mu(q)=\Xi(\mu).
\end{equation}
The same supremum holds over arbitrary finite-dimensional quantum
states and finite-outcome POVMs. For full-support \(\mu\), every
finite-outcome quantum behaviour has statistical strength strictly
smaller than \(\Xi(\mu)\).
\end{theorem}

For uniform inputs, the ceiling is \(\log_2(4/3)\), the relative
entropy of the PR box computed by Grudka et al.\ \cite{Grudka}.
The finite-outcome gap follows from
\(\LF(q)\ge d-\sqrt{d^2-1}>1/(2d)\), a quantitative local-component
bound related to the two-input results of
Gisin, M\'ethot, and Scarani \cite{GMS} and Joshi et al.\ \cite{Joshi}.
For every \(m\ge200\), the explicit comparison
\eqref{eq:KLseparation} shows that uniform-input \(T_m\) has greater
strength than \(Q_{2^m}\) after optimization of the latter's input law.

The projective proof reduces the measurements to fourth roots of the
identity while preserving the original strategy as a reducing
summand. A matrix Hardy inequality then supplies positive forms
indexed by partitions of the circle. A tilted partition law matches
the CGLMP coefficients; polynomial bounds control the remaining
error. Sections~\ref{sec:reduction}--\ref{sec:partition} develop this
argument, Section~\ref{sec:equality} proves rigidity and stability,
and Section~\ref{sec:operational} treats the operational comparisons.
The appendices contain the facet analysis and the certificate data.

\section{Preliminaries and notation}\label{sec:preliminaries}

There are two parties, inputs \(x,y\in\{0,1\}\), and outcomes
\(a,b\in\Z_d\), with \(d\ge2\).
The local polytope \(\Loc_d\) is the convex hull of
\[
 q_s(a,b|x,y)=\ind_{\{a=a_x,\ b=b_y\}},
 \qquad s=(a_0,a_1,b_0,b_1)\in\Z_d^4.
\]
The no-signalling polytope \(\NS_d\) consists of nonnegative normalized
tables whose one-party marginals are independent of the other party's
input.

Write \([u]_d\in\{0,\ldots,d-1\}\) for reduction modulo \(d\), and
orient the four contexts cyclically:
\[
 r_{00}=[a-b]_d,\quad r_{10}=[b-a]_d,\quad
 r_{11}=[a-b]_d,\quad r_{01}=[b-a-1]_d.
\]
Our normalization is
\begin{equation}\label{eq:Id}
 S_d(q)=\sum_{x,y}\mathbb E_qr_{xy},
 \qquad I_d(q)=4-\frac{2S_d(q)}{d-1}.
\end{equation}
For a deterministic strategy, the residuals sum to \(-1\) modulo
\(d\), so \(S_d\ge d-1\), equivalently \(I_d\le2\)
\cite{CGLMP,AGG}.

For the maximally entangled state,
\[
 \ket{\Phi_n}=\frac1{\sqrt n}\sum_{j=0}^{n-1}\ket{jj},
 \qquad \tau_n=\Tr/n,\qquad
 \norm{X}_{2,n}^2=\tau_n(X^\dagger X).
\]
With \(P^0=A_0^T,P^1=B_0,P^2=A_1^T,P^3=B_1\),
\[
 q(a,b|x,y)=\tau_n(P^{2x}_aP^{1+2y}_b).
\]
PVMs may contain zero projections and unequal ranks.
Common conjugation of the four tracial PVMs by \(U\) corresponds
physically to \(\overline U\otimes U\), which preserves \(\Phi_n\).
Operator norms are denoted by \(\norm{\cdot}\), and
\(\norm{\cdot}_F\) is the unnormalized Frobenius norm.
Matrix Hilbert--Schmidt norms and inner products without a dimension
subscript use the normalized trace of the matrix space in question.

The standard tracial PVMs on \(\C^d\) are
\begin{equation}\label{eq:standard}
 \ket{\psi_{t,a}}=\frac1{\sqrt d}\sum_{j=0}^{d-1}
 e^{-2\pi\ii j(a-t/4)/d}\ket j,\qquad
 \Pi^t_a=\ket{\psi_{t,a}}\bra{\psi_{t,a}},
 \quad 0\le t<4.
\end{equation}
Alice's physical bases are conjugated and Bob's are unchanged;
individual basis-vector phases leave the projectors unchanged.
For \(n=dk\), define
\begin{equation}\label{eq:distance}
 \Err_{d,n}(P)=\inf_{U\in U(n)}
 \sum_{t=0}^3\sum_{a=0}^{d-1}
 \norm{P^t_a-U(\Pi^t_a\otimes I_k)U^\dagger}_{2,n}^2.
\end{equation}
The stability constants are
\begin{equation}\label{eq:stabilityconstants}
 \begin{split}
 C_d^{\rm st}
 &=2[1+4(1+2\sqrt d)(1+2\sqrt{d-1})]^2,\\
 R_d^{\rm st}
 &=C_d^{\rm st}
       \left(\sqrt{224}+\frac{8d}{\sqrt{d-1}}\right)^2,\\
 K_d&=\frac34d(d-1)(d+1)R_d^{\rm st}.
 \end{split}
\end{equation}

For the enlarged-space constructions, use
\(\Comp{r}_a=\ket a\bra a\otimes I_r\) on \(\C^d\otimes\C^r\).
The symbols \(E_{j\ell}\) denote spectral projections of fourth-root
unitaries \(V_j=\sum_{\ell=0}^3\ii^\ell E_{j\ell}\).
We reserve \(\mathcal Z\) for the reducing isometry and
\(\Lambda=\mathcal Z\mathcal Z^\dagger\) for its range projection.
The isometry adjoining the spectral-outcome labels is \(\mathcal V\),
with coordinate projections \(\mathsf D_{j\ell}\) on the label
ancilla. The errors used in Section~\ref{sec:equality} are
\[
 \epsanc=\frac1r\sum_{j,\ell}
       \norm{\mathsf D_{j\ell}\mathcal V-\mathcal V E_{j\ell}}_F^2,
 \qquad
 \epsred=\norm{\Lambda-I_d\otimes\bar\Lambda}_{2,N}^2,
\]
where \(\Lambda=\diag(\Lambda_0,\ldots,\Lambda_{d-1})\),
\(\bar\Lambda=d^{-1}\sum_a\Lambda_a\), and \(N=4dn\).
For PVMs \(P,Q\) on \(\C^n\), write
\(\epsmeas(P,Q)=\sum_{t,a}\norm{P^t_a-Q^t_a}_{2,n}^2\).

Let \(u_d(a,b|x,y)=d^{-2}\). Global critical visibility is
\[
 v_{\rm crit}(q)=
 \max\{v\in[0,1]:vq+(1-v)u_d\in\Loc_d\}.
\]
Smaller visibility means greater white-noise resistance.
For the CGLMP functional alone, the corresponding threshold is the
CGLMP critical visibility \(v_{\rm CGLMP}(q)=2/I_d(q)\), the smallest
visibility at which the isotropic mixture still violates the CGLMP
inequality.
For equal-rank PVMs on \(\Phi_n\), this is the visibility of the
isotropic state
\(v\ket{\Phi_n}\bra{\Phi_n}+(1-v)I/n^2\).

An input distribution \(\mu\) may correlate the two parties' inputs,
and is independent of the hidden variable. Statistical strength in
bits per complete trial is
\begin{equation}\label{eq:strength}
 \mathcal S_\mu(q)=\min_{\ell\in\Loc_d}
       \sum_{x,y}\mu_{xy}D(q_{xy}\Vert\ell_{xy}),\qquad
 D(p\Vert q)=\sum_zp_z\log_2(p_z/q_z).
\end{equation}
We use the usual extended-value conventions at zero probabilities.
The local fraction and the input-law function are
\[
 \LF(q)=\max\{\alpha:q\ge\alpha L
       \text{ entrywise for a normalized }L\in\Loc_d\},
\]
\begin{equation}\label{eq:Xi}
 \Xi(\mu)=
 \min_{\substack{0\le s_e\le1\\\sum_es_e=3}}
       -\sum_e\mu_e\log_2s_e,\qquad e\in\{00,01,10,11\}.
\end{equation}
For \(\alpha<1\), the normalized remainder of \(q-\alpha L\) is
no-signalling. We also use
\[
 \TV(p,q)=\frac12\sum_z|p_z-q_z|,\qquad
 H(p,q)^2=\sum_z(\sqrt{p_z}-\sqrt{q_z})^2,
\]
so that \(\TV(p,q)\le H(p,q)\).

\section{Fourth roots and the matrix Hardy inequality}
\label{sec:reduction}

\subsection{The companion representation}

For \(V_0,\ldots,V_{d-1}\in U(r)\) with \(V_j^4=I\), put
\[
 z_{jk}=\tau_r(V_j^\dagger V_k),\qquad
 p_{jk}=\frac{1-\Re z_{jk}+\Im z_{jk}}2,\qquad
 \Fdef=\sum_{j<k}p_{jk}.
\]
For a coefficient sequence \(c=(c_h)_{h=1}^{d-1}\), define
\begin{equation}\label{eq:Lc}
 L_c=\sum_{j<k}
 \bigl[c_{k-j}(1-\Re z_{jk})+c_{d-k+j}\Im z_{jk}\bigr].
\end{equation}
The spectral projections satisfy
\begin{equation}\label{eq:forbidden}
 p_{jk}=\sum_{\ell=0}^3\sum_{u=1}^2
              \tau_r(E_{j\ell}E_{k,\ell+u})\ge0,
\end{equation}
with spectral indices read modulo four.
For \(f(x)=\sec(\pi x/2)\), set
\[
 \begin{gathered}
 P_{d,h}(x)=\binom dhx^h(1-x)^{d-h},\qquad
 B_dg(h)=d\int_0^1P_{d,h}(x)g(x)\dd x,\\
 \beta_h=\frac{f(h/d)}{2d^2},\qquad
 q_h=\frac{B_df(h)}{2d^2},\qquad
 \Gdef=L_\beta,\qquad \Ddef=L_q.
 \end{gathered}
\]

\begin{lemma}[Companion representation]\label{lem:companion}
Every projective strategy in Theorem~\ref{thm:value} has a companion
fourth-root tuple of size \(r=4n\) satisfying
\begin{equation}\label{eq:normalization}
 \IF-I_d=\frac{8d}{d-1}\Gdef.
\end{equation}
The associated physical PVMs act on \(\C^d\otimes\C^{4n}\), and the
original four PVMs form a common reducing summand.
\end{lemma}

\begin{proof}
Let \(J\ket a=\ket{a-1}\) on \(\C^d\), and define
\[
 Z_t=\sum_{c=0}^{d-1}J^c\otimes P^t_c\quad(0\le t<4),
 \qquad Z_4=Z_0(J\otimes I_n).
\]
Each \(Z_t\) is unitary and commutes with \(J\otimes I_n\), and
\[
 Z_t\Comp n_aZ_t^\dagger
 =\sum_{k=0}^{d-1}\ket k\bra k\otimes P^t_{a-k}.
\]
This outcome-covariant strategy preserves the Bell value and contains
the original strategy in its \(k=0\) block.

Put \(\mathsf U_t=Z_t^\dagger Z_{t+1}\). Their product is
\(J\otimes I_n\). On a four-dimensional clock, define
\[
 \mathcal R=\sum_{t=0}^3
 \ket t\bra{t+1\bmod4}\otimes\mathsf U_t.
\]
Every cyclic product of the \(\mathsf U_t\)'s equals \(J\otimes I_n\);
hence \(\mathcal R^4=I_4\otimes J\otimes I_n\).
With
\[
 \mathsf F_{aj}=d^{-1/2}e^{2\pi\ii aj/d},\qquad
 \lambda_j=e^{\pi\ii j/(2d)},
\]
reordering tensor factors and Fourier transforming the outcome
coordinate gives
\[
 \mathcal R=(\mathsf F\otimes I_{4n})
 \diag(\lambda_jV_j)_{j=0}^{d-1}
 (\mathsf F^\dagger\otimes I_{4n}),
 \qquad V_j\in U(4n),\quad V_j^4=I.
\]
In these coordinates the PVMs are
\begin{equation}\label{eq:rootmeasurements}
 \RootP^t_a(V)=\mathcal R^t\Comp{4n}_a\mathcal R^{-t},
 \qquad 0\le t<4.
\end{equation}
Clock averaging identifies their four-edge cost with the sum of the
original transition costs. In clock zero,
\(\mathsf U_0\cdots\mathsf U_{t-1}=Z_0^\dagger Z_t\).
The isometry
\begin{equation}\label{eq:reducingisometry}
 \mathcal Z\xi=\sum_a\ket a\otimes\ket0\otimes P^0_a\xi
\end{equation}
therefore satisfies
\[
 \mathcal Z^\dagger\mathcal Z=I_n,\qquad
 \RootP^t_a(V)\mathcal Z=\mathcal ZP^t_a.
\]
Self-adjointness makes its image a reducing subspace.

The one-edge cost is
\[
 \frac{d-1}{2}
 -\frac1{2d}\sum_{j<k}
 \left[
 \sec\!\left(\frac{\pi(k-j)}{2d}\right)\Re z_{jk}
 -\csc\!\left(\frac{\pi(k-j)}{2d}\right)\Im z_{jk}
 \right],
\]
using
\(\sum_{a=0}^{d-1}a\omega^{ah}=-d/(1-\omega^h)\) for
\(h\not\equiv0\pmod d\).
Subtracting its value at \(V_j=I\) gives \(d\Gdef\).
The four edges and \eqref{eq:Id} yield \eqref{eq:normalization}.
For \(V_j=I\), Fourier conjugation gives \eqref{eq:standard}, and
the same sum evaluates its value as \eqref{eq:sharpvalue}.
\end{proof}

\subsection{Trace-density moments and pinching}

Let \(P_-\) retain the strictly negative Fourier coefficients on
\(\mathbb T\). The \(L^2\)-norm uses normalized Haar measure and
normalized matrix trace.

\begin{theorem}[Fourth-root Hardy bound]\label{thm:hardy}
For every measurable \(U:\mathbb T\to U(r)\) satisfying
\(U(z)^4=z^{-1}I\) almost everywhere,
\begin{equation}\label{eq:Hardy}
 \norm{P_-U}_2^2\ge
 \nu_4:=\frac12-\frac{4\Gc}{\pi^2},\qquad
 \Gc=\sum_{k=0}^\infty\frac{(-1)^k}{(2k+1)^2},
\end{equation}
where \(\Gc\) is Catalan's constant.
\end{theorem}

The moment identity below was obtained independently by
Heinävaara \cite[Lemma 2.5, equation (16)]{Heinavaara}.
We give a proof adapted to matrix pencils.

\begin{lemma}[Trace-density moments]\label{lem:density}
Let \(H=H^\dagger\) be invertible and \(Y=Y^\dagger\) in \(M_r(\C)\).
Counting algebraic multiplicities, define
\[
 w_{H,Y}(y)=\frac1{2\pi r}
 \sum_{\lambda\in\spec(H^{-1}(Y-yI))}|\Im\lambda|.
\]
This function is continuous, nonnegative, and supported on the spectral
interval of \(Y\). For every integer \(m\ge0\),
\begin{equation}\label{eq:densitymoments}
 \int y^m w_{H,Y}(y)\dd y
 =\frac1{(m+1)(m+2)!}
 \left.\frac{\mathrm d^{m+2}}{\mathrm dt^{m+2}}
           \tau_r(H+tY)_+^{m+1}\right|_{t=0}.
\end{equation}
\end{lemma}

\begin{proof}
For \(\Im z>0\), an eigenvector of \(L(z)=H^{-1}(Y-zI)\) satisfies
\[
 -\Im z\,\norm v^2=\Im\lambda\,v^\dagger Hv.
\]
Its spectrum therefore avoids the real axis. The sum
\[
 g(z)=r^{-1}\sum_{\Im\lambda(L(z))>0}\lambda(L(z))
\]
is holomorphic, being the trace of \(L(z)\) times its upper-cluster
Riesz projection, with algebraic multiplicities also at repeated or
defective eigenvalues.

The sum extends continuously to the real axis. Choose disjoint small
disks around the limiting eigenvalues. In a connected upper
half-neighbourhood, the number of selected roots in each disk is
constant. A real limiting cluster contributes an integer multiple
of its real eigenvalue. At real \(y\), the characteristic polynomial
\(\det(Y-yI-\lambda H)/\det H\) has real coefficients, so nonreal
roots occur in conjugate pairs. Consequently
\[
 \Im g(y+\ii0)=\pi w_{H,Y}(y).
\]
Outside the spectral interval of \(Y\), definiteness of \(Y-yI\)
makes the pencil similar, up to sign, to a Hermitian matrix, and
the density vanishes.

Near \(w=0\), put
\[
 K(w)=H^{-1}(I-wY),\qquad
 q(w)=\tau_r(K(w)\Pi_-(w)),
\]
where \(\Pi_-(w)\) continues the negative spectral cluster of \(H^{-1}\).
For small real \(w\), similarity to
\((I-wY)^{1/2}H^{-1}(I-wY)^{1/2}\) shows that \(q\) has real Taylor
coefficients. In the upper exterior domain, \(g(z)=-zq(1/z)\).
Writing \(q(w)=\sum_jq_jw^j\), Schwarz reflection and Cauchy's formula
around the spectral slit give
\begin{equation}\label{eq:Cauchydensity}
 g(z)=-q_0z-q_1+\int\frac{w_{H,Y}(y)}{y-z}\dd y,\qquad
 \int y^m w_{H,Y}(y)\dd y=[w^{m+2}]q(w).
\end{equation}
The endpoint arcs vanish because pencil eigenvalues remain bounded
near finite points; continuity controls the two banks of the slit.

Let \(k=m+2\) and \(\mathsf R_H(\lambda)=(\lambda I-H)^{-1}\).
Contours enclose the indicated positive or negative spectra and
exclude zero. Expanding the logarithmic determinant ratio gives
\begin{align*}
 [w^k]q(w)
 &=\frac{(-1)^k}{kr\,2\pi\ii}
   \oint_{\Gamma_{H^{-1},-}}
   \Tr[((\zeta H-I)^{-1}Y)^k]\dd\zeta\\
 &=-\frac1{kr\,2\pi\ii}
   \oint_{\Gamma_{H,-}}
   \lambda^{k-2}\Tr[(\mathsf R_H(\lambda)Y)^k]\dd\lambda\\
 &=\frac1{kr\,2\pi\ii}
   \oint_{\Gamma_{H,+}}
   \lambda^{k-2}\Tr[(\mathsf R_H(\lambda)Y)^k]\dd\lambda.
\end{align*}
The second line substitutes \(\lambda=1/\zeta\); the third uses
decay \(O(\lambda^{-2})\) at infinity.
Positive-cluster functional calculus and integration by parts, with
\[
 \frac{\mathrm d}{\mathrm d\lambda}
 \Tr[(\mathsf R_H(\lambda)Y)^k]
 =-k\Tr[(\mathsf R_H(\lambda)Y)^k\mathsf R_H(\lambda)],
\]
give
\[
 [t^k]\tau_r(H+tY)_+^{k-1}=(k-1)[w^k]q(w).
\]
Together with \eqref{eq:Cauchydensity}, this proves
\eqref{eq:densitymoments}.

For the two-node specialization, expansion at \(a\) gives
\begin{equation}\label{eq:twonoderesidue}
 \underset{\lambda=a}{\operatorname{Res}}\,
 \frac{\lambda^{m-2}}{(\lambda-a)^p(\lambda-b)^q}
 =
 \binom{m-2}{p-1}
 \frac{a^{q-1}(-b)^{p-1}}{(a-b)^{m-1}},
 \qquad p,q\ge1,\quad p+q=m,\quad a>0>b.
\end{equation}
\end{proof}

\begin{lemma}[Positive pinching identity]\label{lem:pinching}
Let \(B=B^\dagger\), \(B^2=I\), \(Y=Y^\dagger\), and
\[
 C=\Pi_B(Y):=\frac{Y+BYB}{2},\qquad D=Y-C.
\]
There is a nonnegative finite measure \(\nu\), supported on
\([-1,1]\) times the spectral interval of \(Y\), such that
\begin{equation}\label{eq:pinching}
 \tau_r f(sB+tY)-\tau_r f(sB+tC)
 =t^2\int f''(-sx+ty)\dd\nu(x,y)
\end{equation}
for polynomials \(f\). Its density on \(|x|<1\) is
\begin{equation}\label{eq:nudensity}
 \frac{\omega(x,y)}{1-x^2},\qquad
 \omega(x,y)=\frac1{2\pi r}
 \sum_{\lambda\in\spec((B-xI)(Y-yI))}|\Im\lambda|,
\end{equation}
and
\begin{equation}\label{eq:numajorant}
 0\le\frac{\omega(x,y)}{1-x^2}
 \le\frac{\norm D}{2\pi\sqrt{1-x^2}},\qquad
 \nu(\R^2)=\frac12\tau_r(D^2).
\end{equation}
\end{lemma}

\begin{proof}
Apply Lemma~\ref{lem:density} with \(H=(B-xI)^{-1}\).
In a basis diagonalizing \(B\), let \(p\) entries of an index cycle
\(\mathbf i=(i_1,\ldots,i_m)\) belong to its \(+1\) eigenspace and
let \(q=m-p\). Put
\(Y_{\mathbf i}=Y_{i_1i_2}\cdots Y_{i_mi_1}\).
Substituting \(a=(1-x)^{-1}\), \(b=-(1+x)^{-1}\) in
\eqref{eq:twonoderesidue} yields
\begin{equation}\label{eq:twonodemoments}
 \int y^{m-2}\frac{\omega(x,y)}{1-x^2}\dd y
 =\frac1{rm\,2^{m-1}}
 \sum_{\mathbf i:\,1\le p<m}
 \binom{m-2}{p-1}
 (1-x)^{p-1}(1+x)^{q-1}Y_{\mathbf i}.
\end{equation}

The coefficient of \(t^m\) in the mixed-cycle part of
\(\tau_r e^{sB+tY}\) is
\[
 \frac1{rm}\sum_{\mathbf i:\,1\le p<m}Y_{\mathbf i}
 \int_{\substack{u_j\ge0\\\sum_ju_j=1}}
 e^{s\sum_jb_{i_j}u_j}\dd u.
\]
Grouping the positive-\(B\) coordinates with sum \(v\) gives a beta
integral with density
\(v^{p-1}(1-v)^{q-1}/((p-1)!(q-1)!)\).
The substitution \(x=1-2v\), followed by
\eqref{eq:twonodemoments}, gives the coefficient on the right of
\eqref{eq:pinching}, including the sign in \(-sx\).
Comparison of coefficients proves the polynomial identity.

A block-diagonal similarity transforms the pencil in
\eqref{eq:nudensity} into a matrix with Hermitian diagonal blocks
and off-diagonal blocks
\[
 \sqrt{1-x^2}D_{+-},\qquad
 -\sqrt{1-x^2}D_{+-}^\dagger.
\]
Its imaginary part has norm at most
\(\sqrt{1-x^2}\norm D\). The same bound holds for the imaginary part
of each eigenvalue, proving the majorant in \eqref{eq:numajorant}.
Integration of the \(m=2\) case of \eqref{eq:twonodemoments} in \(x\)
gives the total mass.
\end{proof}

\begin{corollary}[Holomorphic pinching on a strip]\label{cor:holomorphic}
Under the hypotheses of Lemma~\ref{lem:pinching},
identity \eqref{eq:pinching} holds with \(s=\rho\in[0,1)\) and
\(t=\ii\) for every \(f\) holomorphic on \(|\Re z|<1\).
\end{corollary}

\begin{proof}
Choose \(\rho<\rho'<1\) and \(L>\norm Y\).
The rectangle
\(\{z:|\Re z|\le\rho',\ |\Im z|\le L\}\)
contains the spectra of \(\rho B+\ii Y,\rho B+\ii C\) and all scalar
arguments \(-\rho x+\ii y\) in the support of \(\nu\).
Sum the polynomial identity for a large-parameter resolvent, continue
in the connected exterior of this rectangle, and apply Cauchy's
formula on its boundary.
\end{proof}

\subsection{Proof of the Hardy bound}

\begin{proof}[Proof of Theorem~\ref{thm:hardy}]
On \(|\Re z|<1\), define \(h\) by
\[
 h(0)=\frac{8\Gc}{\pi^2},\qquad h'(0)=0,\qquad
 h''(z)=-\sec(\pi z/2).
\]
Its real part \(u(x,y)=\Re h(x+\ii y)\) extends continuously to the
closed strip and satisfies
\begin{equation}\label{eq:stripscalar}
 u(\pm1,y)=|y|,\qquad
 0\le u(x,y)\le\frac{8\Gc}{\pi^2}+|y|.
\end{equation}
Indeed,
\[
 h'(z)=-\frac2\pi
 \log\bigl(\sec(\pi z/2)+\tan(\pi z/2)\bigr).
\]
The boundary imaginary parts give the slopes of \(|y|\), and
\(\int_0^1(1-x)\sec(\pi x/2)\dd x=8\Gc/\pi^2\) fixes the constant.
Also \(u_{xx}=-\Re\sec(\pi(x+\ii y)/2)<0\).
Symmetry in \(x\) and integration of \(\sech(\pi y/2)\) give
\eqref{eq:stripscalar}.

For a matrix \(X\) with spectrum in the closed strip, put
\[
 \mathscr U(X)=r^{-1}\sum_{\lambda\in\spec X}\Re h(\lambda).
\]
Apply Corollary~\ref{cor:holomorphic} and let \(\rho\uparrow1\).
Since \(\Pi_B(Y)\) commutes with \(B\), the boundary values in
\eqref{eq:stripscalar} give
\begin{equation}\label{eq:stripmatrix}
 \mathscr U(B+\ii Y)-\tau_r|\Pi_B(Y)|
 =
 \int\Re\sec\!\left(\frac{\pi(x+\ii y)}2\right)\dd\nu(x,y)
 \ge0.
\end{equation}
For the limit, combine \eqref{eq:numajorant} with
\[
 \int_\R\Re\sec\!\left(\frac{\pi(\rho x+\ii y)}2\right)\dd y=2.
\]
The absolute integral of each \(x\)-slice is at most
\(\norm D/(\pi\sqrt{1-x^2})\).
Restriction to smaller strips followed by passage to the endpoints
proves convergence for the possibly nonnormal matrix \(B+\ii Y\).

Set \(\mathcal W(s)=sU(s^4)\) on the circle and
\[
 A=\Re\mathcal W+\Im\mathcal W,\qquad
 B=\Re\mathcal W-\Im\mathcal W.
\]
These are commuting Hermitian involutions.
The identity \(\mathcal W(\ii s)=\ii\mathcal W(s)\) makes their means
zero. For the circular Hilbert transform with multiplier
\(-\ii\operatorname{sgn}(k)\), Parseval gives
\begin{equation}\label{eq:HardyParseval}
 1-2\norm{P_-U}_2^2
 =\int_{\mathbb T}\tau_r(A\,\Hilb B)\dd m
 \le\int_{\mathbb T}\tau_r|\Pi_B(\Hilb B)|\dd m.
\end{equation}
The frequencies of \(\mathcal W\) are \(4k+1\), with the same
positive/negative split as those of \(U\). Expanding
\(\mathcal W=((1+\ii)A+(1-\ii)B)/2\) proves the equality.

The analytic matrix function with boundary values \(B+\ii\Hilb B\)
has value zero at the disk centre. Trace holomorphic functional
calculus and the mean-value identity give mean
\(\mathscr U=8\Gc/\pi^2\).
For boundary passage, use radial limits and
\[
 0\le\mathscr U(X)\le\frac{8\Gc}{\pi^2}+\tau_r|\Im X|.
\]
This follows from \eqref{eq:stripscalar}, Schur triangularization,
and convexity of absolute value on the diagonal.
The radial Hilbert transforms are uniformly \(L^2\)-bounded, hence
the displayed quantities are uniformly integrable.
Applying \eqref{eq:stripmatrix} in \eqref{eq:HardyParseval} gives
\(1-2\norm{P_-U}_2^2\le8\Gc/\pi^2\), proving \eqref{eq:Hardy}.
\end{proof}

\section{Positive partition forms and the CGLMP bound}
\label{sec:partition}

\begin{theorem}[Fourth-root bound]\label{thm:certificate}
Let \(d\ge2\), \(r\ge1\), and
\(V_0,\ldots,V_{d-1}\in U(r)\) satisfy \(V_j^4=I\).
Then
\begin{equation}\label{eq:strongcertificate}
 \Gdef\ge\left(\frac32-\frac\pi4\right)\Ddef
       +\frac{\Fdef}{6d^2(d+1)},\qquad \Ddef\ge0.
\end{equation}
\end{theorem}

\subsection{From Hardy space to circular partitions}

\begin{lemma}[Partition forms]\label{lem:partitions}
For a consecutive partition
\(\ell=(I_0,\ldots,I_{d-1})\) of \([0,1]\), the fourth-root tuple
defines a nonnegative form
\begin{equation}\label{eq:partitionform}
 \begin{split}
 \mathcal K_\ell(V)
 &:=
 \norm{P_-[e^{-\pi\ii t/2}V(t)^\dagger]}_2^2-\nu_4\\
 &=\frac12\sum_{j<k}\int_{I_j}\int_{I_k}
 \left[f(t-s)(1-\Re z_{jk})
       +f(1-t+s)\Im z_{jk}\right]\dd t\dd s\ge0,
 \end{split}
\end{equation}
where \(V(t)=V_j\) on \(I_j\).
Extend \(V_{j+d}=-\ii V_j\). Averaging over the \(d\) cyclic
assignments to a circular partition gives \(L_c\), with
\begin{equation}\label{eq:arccount}
 c_h=\frac1{2d}\int_0^1f(x)\Pr(N_x=h)\dd x,\qquad 1\le h<d,
\end{equation}
where \(N_x\) counts boundaries in a length-\(x\) arc from a uniform
origin. In particular, \(\Ddef\ge0\).
\end{lemma}

\begin{proof}
The function \(U(e^{2\pi\ii t})=e^{-\pi\ii t/2}V(t)^\dagger\)
satisfies \(U(z)^4=z^{-1}I\), so \eqref{eq:Hardy} gives positivity.
For \(s<t\), put \(x=t-s\) and use
\[
 e^{\pi\ii x/2}(1+\ii\cot\pi x)
 =\frac{f(x)+\ii f(1-x)}2.
\]
Abel regularization and pairing the two half-squares in
\(P_-=(I-P_0-\ii\Hilb)/2\) give \eqref{eq:partitionform}, with constant
\(\frac12-\frac12\int_0^1(1-x)f(x)\dd x=\nu_4\).
The integrals are finite: near the diagonal the imaginary part
vanishes within an interval, while cross-interval regions have length
\(O(x)\). This also justifies the regularization limit.

The cyclic average uses the zero Hardy gap of scalar cuts.
For the periodic branch \(U_0(t)=e^{-\pi\ii t/2}\), \(0\le t<1\),
\[
 \norm{P_-U_0}_2^2
 =\frac8{\pi^2}\sum_{n\ge0}\frac1{(4n+3)^2}
 =\nu_4.
\]
A scalar label dropping by one power of \(\ii\) at \(T\) gives
\begin{equation}\label{eq:rotatedbranch}
 e^{-\pi\ii t/2}\ii^{-a+\ind_{\{t>T\}}}
 =
 \ii^{1-a}e^{-\pi\ii T/2}U_0([t-T]_1)
 \quad\text{almost everywhere}.
\end{equation}
Translation and multiplication by a constant phase preserve Fourier
moduli, so \eqref{eq:rotatedbranch} has the same negative Fourier
energy.

For a pair of intervals separated by \(h\) boundaries, the direct arc
contributes \(f(x)\) and the complementary arc contributes \(f(1-x)\).
Crossing the label origin replaces \(z\) by \(-\ii\bar z\), interchanging
these coefficients. Subtraction of the value at \(V_j=I\) removes
the scalar-cut contribution, and averaging the starting label gives
the factor \(1/d\). This proves \eqref{eq:arccount}.
Independent uniform boundaries have
\(\Pr(N_x=h)=P_{d,h}(x)\), giving \(L_c=\Ddef\ge0\).
\end{proof}

\subsection{The tilted law and coefficient matching}

Define
\begin{equation}\label{eq:vkernel}
 v(x)=-\frac{x(1-x)}2+
 \frac{2x-1}{\pi[\sec(\pi x/2)-\csc(\pi x/2)]},
\end{equation}
with removable singularities filled continuously.
Also put
\[
 b(x)=-\frac12\left[x(1-x)f(x)-\frac{2x}{\pi}\right]'',
 \qquad
 S(x)=\frac{1-z\cot z}{\sin^2z},\quad
 z=\frac\pi2(x-\tfrac12),\quad S(1/2)=\frac13,
\]
and define
\begin{equation}\label{eq:es}
 e_{d,h}=f(h/d)-B_df(h)-\frac1dB_db(h),\qquad
 s_{d,h}=\frac1dB_dS(h).
\end{equation}

\begin{lemma}[Tilted partition law]\label{lem:tilted}
There is a probability density on circular partitions whose averaged
Hardy form satisfies
\begin{equation}\label{eq:Kv}
 \mathcal K_d^v
 =L_{\{B_df+(1/d)B_d(vf)''\}/(2d^2)}
 \ge\cstar\Ddef,\qquad \cstar=\frac32-\frac\pi4.
\end{equation}
Moreover,
\begin{equation}\label{eq:matching}
 \Gdef=\mathcal K_d^v+\frac1{2d^2}L_{e_d+s_d},
 \qquad s_{d,h}=s_{d,d-h}\ge\frac1{3(d+1)}.
\end{equation}
\end{lemma}

\begin{proof}
The function \(v\) is symmetric about \(1/2\), has zero value and
first derivative at both endpoints, and
\(v''(0)=v''(1)=(\pi-2)/2\).
It has the expansion
\begin{equation}\label{eq:vFourier}
 v(x)=\sum_{n\ge1}a_n\frac{\sin^2(\pi nx)}{\pi^2n^2},
 \qquad a_n>0,\qquad
 \sum_{n\ge1}a_n=\frac\pi4-\frac12.
\end{equation}
To prove this, write \(u=x-\tfrac12\), so that
\[
 v(x)=\frac{u^2-\tfrac14}{2}
 +\frac{u}{\sqrt2\pi}
 \left[\csc(\pi u/2)-2\sin(\pi u/2)\right].
\]
The nonconstant even coefficients of \(u\csc(\pi u/2)\) are positive.
For \(a=\pi/2\),
\[
 -[u\sin(au)]^{(4)}
 =a^3[4\cos(au)-au\sin(au)]>0
 \quad(|au|\le\pi/4).
\]
Thus \(v^{(4)}>0\), and four integrations by parts give
\[
 a_n=\frac1{4\pi^2n^2}
 \int_0^1v^{(4)}(x)[1-\cos(2\pi nx)]\dd x>0.
\]
The coefficients are \(O(n^{-2})\).
Fourier inversion and evaluation of \(v''(0)\) prove
\eqref{eq:vFourier}.

Tilt independent uniform boundary points \(T_1,\ldots,T_d\) by
\begin{equation}\label{eq:partitiondensity}
 \varpi_d=1+\frac1d\sum_{a<b}v''(T_a-T_b)
 =\cstar+\frac1d\sum_{n\ge1}a_n
       \left|\sum_{j=1}^de^{2\pi\ii nT_j}\right|^2,
\end{equation}
using the periodic extension in the differences.
Its expectation is one and \(\varpi_d\ge\cstar>0\).
For one Fourier-power density, the count generating function is
\[
 (1-x+xz)^d+
 (d-1)\frac{\sin^2(\pi nx)}{\pi^2n^2}
       (z-1)^2(1-x+xz)^{d-2}.
\]
Summing with \eqref{eq:vFourier} shows that under
\eqref{eq:partitiondensity},
\begin{equation}\label{eq:tiltedcount}
 \Pr_{\varpi_d}(N_x=h)=
 P_{d,h}(x)+\frac{v(x)}dP_{d,h}''(x).
\end{equation}
Equations \eqref{eq:arccount} and \eqref{eq:tiltedcount}, followed by
two integrations by parts, give \eqref{eq:Kv}.
The endpoint terms vanish because \(vf\) vanishes at both endpoints,
\((vf)'\) has finite one-sided limits, and \(P_{d,h}\) vanishes there.

Differentiating \eqref{eq:vkernel} gives \((vf)''=b-S\).
The Taylor series of \(z\cot z\) gives
\(S(x)=S(1-x)\ge1/3\). Substitution in \eqref{eq:es} proves
\eqref{eq:matching}.
\end{proof}

\subsection{Circulation and the scalar estimate}

\begin{lemma}[Circulation bound]\label{lem:circulation}
For \(i<j<k\), let
\(C_{ijk}=\Im z_{ij}+\Im z_{jk}-\Im z_{ik}\).
Then
\begin{equation}\label{eq:circulation}
 |C_{ijk}|\le5(p_{ij}+p_{jk}+p_{ik}).
\end{equation}
For \(t_h=e_{d,d-h}-e_{d,h}\) and
\(\omega_{ijk}=t_{j-i}+t_{k-j}-t_{k-i}\),
\begin{equation}\label{eq:Hodge}
 L_e=2\sum_{j<k}e_{d,k-j}p_{jk}
       +\frac1d\sum_{i<j<k}\omega_{ijk}C_{ijk}.
\end{equation}
\end{lemma}

\begin{proof}
The sum-of-squares identity \eqref{eq:circulationSOS} in
Appendix~\ref{app:circulation} proves \eqref{eq:circulation}.
The skew edge array \(t\) has zero divergence.
Expanding \eqref{eq:Lc} over triangles then gives \eqref{eq:Hodge}.
\end{proof}

The absolute triangle load at separation \(h\) is
\begin{equation}\label{eq:load}
 \ell_{d,h}=
 \sum_{k=1}^{h-1}|t_k+t_{h-k}-t_h|
 +\sum_{k=h+1}^{d-1}|t_h+t_{k-h}-t_k|.
\end{equation}
For \(b_d=\max_h|e_{d,h}|\), the bounds
\(|\omega_{ijk}|\le6b_d\) and \eqref{eq:circulation} give
\begin{equation}\label{eq:coarseerror}
 L_e\ge-32b_d\Fdef,
\end{equation}
since each edge belongs to \(d-2\) triangles.

\begin{proof}[Proof of Theorem~\ref{thm:certificate}]
Lemma~\ref{lem:partitions} gives \(\Ddef\ge0\).
Use the normalized scalar functions \(\mathcal E_\theta\) in
Appendix~\ref{app:continuous}, whose normalization
\eqref{eq:Enormalization} and bound \eqref{eq:continuouslemma}
hold on the indicated domains.
For \(d\ge3\), take
\[
 \theta=1/d,\qquad
 (x_1,x_2,x_3)=
 \left(\frac{j-i}{d},\frac{k-j}{d},\frac{d-k+i}{d}\right).
\]
Then
\[
 \Omega_\theta=-\frac{(d+1)(d-2)}d\,\omega_{ijk}.
\]
Put \(b_d^\triangle=\max_{i<j<k}|\omega_{ijk}|\).
The independent choice of \(y\) in \eqref{eq:continuouslemma} gives,
for every \(h\),
\[
 \frac1{3(d+1)}+e_{d,h}
 -\frac{5(d-2)}{2d}b_d^\triangle
 \ge\frac1{6(d+1)}.
\]
Applying Lemma~\ref{lem:circulation} and counting \(d-2\) triangles
per edge yields
\[
 L_{e+s}\ge\frac{\Fdef}{3(d+1)}.
\]
For \(d=2\), the triangle sum is empty and
\eqref{eq:rectanglebound} gives the same conclusion.
Lemma~\ref{lem:tilted}, through \eqref{eq:Kv} and
\eqref{eq:matching}, now proves \eqref{eq:strongcertificate}.
\end{proof}

\begin{remark}[Alternative coefficient]\label{rem:alternative}
A beta--Stein expansion for \(d\ge66\), together with an interval
calculation for \(2\le d\le65\), gives
\begin{equation}\label{eq:piecewisecertificate}
 \Gdef\ge\cstar\Ddef+\eta_d\Fdef,\qquad
 \eta_d=
 \begin{cases}
 [6d^2(d+1)]^{-1},&2\le d\le65,\\
 (2d-131)[6d^2(d+1)(d+2)]^{-1},&d\ge66.
 \end{cases}
\end{equation}
In particular,
\(\eta_d\ge[6d^2(d+1)(d+2)]^{-1}\).
The scalar estimates and derivation are given in
Appendix~\ref{app:scalar}.
\end{remark}

\section{Rigidity and same-dimension stability}\label{sec:equality}

\subsection{Equality and multiplicity}

\begin{valueagain}
Let \(d\ge2\), \(n\ge1\), and let the four measurements on \(\Phi_n\)
be \(d\)-outcome PVMs. Then \(I_d\le\IF\), with \(\IF\) given by
\eqref{eq:sharpvalue}. Equality holds exactly when \(n=dk\) and the
tracial PVMs are commonly unitarily conjugate to
\(\Pi^t_a\otimes I_k\). All projective maximizers on \(\Phi_d\) are
rank-one Fourier-phase measurements.
\end{valueagain}

\begin{proof}[Proof of Theorem~\ref{thm:value}]
Lemma~\ref{lem:companion} and Theorem~\ref{thm:certificate} give the
bound. Equality forces \(\Fdef=0\), so all forbidden products in
\eqref{eq:forbidden} vanish.

For two spectral PVMs \(E_a,F_b\), the allowed columns of row \(a\)
are \(a,a-1\) modulo four. Distinct rows share at most one allowed
column. For \(a\ne c\),
\[
 0=E_aE_c=\sum_bE_aF_bE_c
\]
has at most one potentially nonzero term, which must vanish.
Thus the two PVMs commute. Applying this argument to every pair
gives simultaneous commutativity.

On each joint scalar eigenspace, the labels change at most once,
by decreasing the phase by \(\pi/2\). Hence they are cuts:
\begin{equation}\label{eq:cuts}
 V_j=\ii^{a-\ind_{\{j>j_0\}}},
 \qquad a\in\Z_4,\qquad 0\le j_0<d.
\end{equation}
For \(s=d-1-j_0\),
\[
 \lambda_j\ii^{-\ind_{\{j>j_0\}}}
 =\lambda_s^{-1}\lambda_{j+s\bmod d}.
\]
This frequency shift is implemented by a diagonal unitary in the
outcome coordinate. Every cut therefore gives the standard PVMs
and zero gap, and a direct sum of cuts gives their standard
amplification.

The algebra generated by the standard PVMs is \(M_d(\C)\).
In the gauge where the first PVM is computational, the next basis
is \(U_{\rm F}=\mathsf F\diag(\lambda_j)\mathsf F^\dagger\), with
\[
 (U_{\rm F})_{ab}
 =\frac{1-\ii}{d[1-e^{2\pi\ii(a-b+1/4)/d}]}.
\]
Every entry is nonzero.
An operator commuting with the first PVM is diagonal; commuting
with the rank-one projectors of this second basis forces its
diagonal entries to agree. The commutant of an amplification is
therefore \(I_d\otimes M_k(\C)\). Its reducing projections have ranks
divisible by \(d\), and every restriction is a standard amplification.
Apply this to the reducing isometry \eqref{eq:reducingisometry}.
It gives \(d\mid n\) and the classification in the original dimension.
The converse follows from \eqref{eq:standard}.
\end{proof}

\subsection{Rounding spectral measurements}

The constants in \eqref{eq:stabilityconstants} arise in two stages.
Rounding and descent give \(C_d^{\rm st}\) and \(R_d^{\rm st}\) in
terms of \(\Fdef\). Conversion to a Bell-value deficit uses the
coefficient \(1/[6d^2(d+1)]\) in \eqref{eq:strongcertificate}, and
supplies \(K_d\), the arithmetic bound, the rank-forcing threshold,
and \(\widehat K_d\) in \eqref{eq:KLtradeoff}.

\begin{lemma}[Rounding to cuts]\label{lem:rounding}
For a fourth-root tuple as in Theorem~\ref{thm:certificate}, the
squared distance, summed over sites and spectral outcomes, from
\((E_{j\ell})\) to a commuting cut tuple on the same \(\C^r\) is at
most \(C_d^{\rm st}\Fdef\).
For the corresponding fourth-root unitaries \(D_j\),
\[
 \sum_j\norm{V_j-D_j}_{2,r}^2\le2C_d^{\rm st}\Fdef.
\]
\end{lemma}

\begin{proof}
For two \(q\)-outcome PVMs whose distinct allowed rows have at most
one common allowed column, total forbidden overlap \(f_{\rm bad}\)
bounds
\[
 \sum_{\ell,m}\norm{[E_\ell,F_m]}_2^2\le4qf_{\rm bad}.
\]
For \(\ell\ne\ell'\), put \(x_m=E_\ell F_mE_{\ell'}\).
The identity \(\sum_mx_m=0\) has at most one term with both endpoints
allowed. Bound that term by Cauchy--Schwarz in the sum of the others,
and sum over \(\ell,\ell',m\), using
\[
 \sum_{\ell,m}\norm{[E_\ell,F_m]}_2^2
 =2\sum_{\ell\ne\ell',m}\norm{E_\ell F_mE_{\ell'}}_2^2.
\]
For the present graph this gives
\begin{equation}\label{eq:commutators}
 \Delta_{\rm com}:=
 \sum_{j<k,\ell,m}\norm{[E_{j\ell},E_{km}]}_2^2
 \le16\Fdef.
\end{equation}

For \(m\) PVMs, coherently adjoin the successive outcome labels:
\[
 \mathcal V\xi=
 \sum_{\boldsymbol\ell}
 E_{m-1,\ell_{m-1}}\cdots E_{0,\ell_0}\xi
 \otimes\ket{\boldsymbol\ell}.
\]
This is an isometry. For the coordinate projections
\(\mathsf D_{j\ell}\), the error defined in the preliminaries obeys
\begin{equation}\label{eq:ancillaerror}
 \epsanc\le2\Delta_{\rm com}.
\end{equation}
To prove this estimate, let \(T_i\) be orthogonal projections,
\(x_i=T_ix_{i-1}\), and set
\[
 \epsilon=\sum_i\norm{(I-T_i)x}^2,\qquad
 L=\norm x^2-\norm{x_m}^2,\qquad
 D=\norm x^2-\Re\ip{x}{x_m}.
\]
Telescoping and Cauchy--Schwarz give
\(D\le\sqrt{\epsilon L}\), while \(L\le2D\), hence \(D\le2\epsilon\).
Apply this on Hilbert--Schmidt space to the pinching maps.
The identity
\[
 \norm{X-\sum_\ell E_\ell XE_\ell}_2^2
 =\frac12\sum_\ell\norm{[E_\ell,X]}_2^2
\]
gives \eqref{eq:ancillaerror}.

Pinch \(\mathcal V\mathcal V^\dagger\) by the joint ancilla
projections, obtaining a positive contraction \(\Gamma\) of trace
\(r\), with \(r^{-1}\Tr(\Gamma-\Gamma^2)\le\epsanc\).
Within the ancilla blocks, choose the projection \(\Gamma'\) onto
the \(r\) largest eigenvalues. If
\[
 \zeta_{\rm tail}=\sum_{i>r}\lambda_i
                 =\sum_{i\le r}(1-\lambda_i),
\]
then
\[
 \zeta_{\rm tail}\le\sum_i\lambda_i(1-\lambda_i),\qquad
 \norm{\mathcal V\mathcal V^\dagger-\Gamma'}_F^2
 =2\zeta_{\rm tail}.
\]
A closest isometry \(\mathcal T\) onto \(\operatorname{ran}\Gamma'\)
satisfies \(r^{-1}\norm{\mathcal T-\mathcal V}_F^2\le2\epsanc\).
Compression of the ancilla projections gives commuting PVMs
\(Q_{j\ell}\) on the original space, satisfying
\begin{equation}\label{eq:commutinground}
 \sum_{j,\ell}\norm{E_{j\ell}-Q_{j\ell}}_2^2
 \le2(1+2\sqrt m)^2\Delta_{\rm com}.
\end{equation}
Indeed, the square root of the left side is at most
\(\sqrt{\epsanc}+\sqrt{2m/r}\norm{\mathcal T-\mathcal V}_F\).
Zero singular values of the polar overlap are handled by completing
the partial isometry between the remaining equal-dimensional spaces.

Take \(m=d\), and call the left side of \eqref{eq:commutinground}
\(g\). Completeness and the triangle inequality for concatenated
forbidden products give
\[
 \sqrt{\mathfrak F(Q)}
 \le\sqrt{\mathfrak F(E)}+\sqrt{2(d-1)g}.
\]
On a joint eigenspace of the commuting PVMs, let \(b\) count forbidden
pairs in the label string. Delete their left endpoints, at most
\(b\) positions. The retained labels extend to a cut.
Changing one label costs squared PVM distance two, so the repair
cost is at most \(2\mathfrak F(Q)\).
Combining with \eqref{eq:commutators} gives
\(C_d^{\rm st}\Fdef\). Replacing labels by fourth roots multiplies
the squared bound by at most two.
\end{proof}

\subsection{Descent and the arithmetic gap}

\begin{robustagain}
For \(d\)-outcome projective measurements on \(\Phi_n\), put
\(\Delta=\IF-I_d\) and \(s=n\bmod d\).
If \(d\mid n\), then
\[
 \Err_{d,n}(P)\le\min\{8,K_d\Delta\},
 \qquad K_d\le786432d^6.
\]
For arbitrary \(n\),
\[
 \Delta\ge\frac{s(d-s)}{3n\,d^4(d+1)C_d^{\rm st}},
\]
and, for \(s>0\), the projective maximum satisfies
\[
 \frac{s(d-s)}{3n\,d^4(d+1)C_d^{\rm st}}
 \le\IF-\mathcal I_{d,n}
 \le\frac{s}{n}(\IF-2).
\]
\end{robustagain}

\begin{proof}[Proof of Theorem~\ref{thm:robust}]
Equations \eqref{eq:strongcertificate} and \eqref{eq:normalization}
give, for the companion,
\begin{equation}\label{eq:compdefect}
 \mathfrak F_{\rm comp}
 \le\frac34d(d-1)(d+1)\Delta.
\end{equation}
Fourier unitarity and Lemma~\ref{lem:rounding} imply
\[
 \norm{\mathcal R(V)-\mathcal R(D)}_{2,dr}^2
 \le\frac{2C_d^{\rm st}}d\Fdef.
\]
For unitaries \(U_1,U_2\),
\[
 \sum_a\norm{U_1\Comp r_aU_1^\dagger
             -U_2\Comp r_aU_2^\dagger}_2^2
 \le2\norm{U_1-U_2}_2^2.
\]
Telescoping powers in \eqref{eq:rootmeasurements} gives
\begin{equation}\label{eq:physicalrounding}
 \begin{split}
 \sum_a\norm{\RootP^t_a(V)-\RootP^t_a(D)}_{2,dr}^2
 &\le\frac{4t^2C_d^{\rm st}}d\Fdef,\\
 \sum_{t,a}\norm{\RootP^t_a(V)-\RootP^t_a(D)}_{2,dr}^2
 &\le\frac{56C_d^{\rm st}}d\Fdef.
 \end{split}
\end{equation}

Standardize the cut benchmark by the controlled diagonal unitary
following \eqref{eq:cuts}. It commutes with the computational PVM.
Expressing \(\mathcal Z\) in this gauge, its range projection is
\[
 \Lambda=\mathcal Z\mathcal Z^\dagger
 =\diag(\Lambda_0,\ldots,\Lambda_{d-1}),\qquad
 \rank\Lambda=n,\qquad N=4dn.
\]
Use \(\bar\Lambda=d^{-1}\sum_a\Lambda_a\) and \(\epsred\) as defined
in the preliminaries. For the second standard basis, let
\(B_{ab}=|(U_{\rm F})_{ab}|^2\). Its Fourier eigenvalues are
\[
 \widehat B_k=e^{\pi\ii k/(2d)}
 \left(1-\frac kd-\ii\frac kd\right),
\]
so \(BB^T\) has eigenvalues \(1-2k(d-k)/d^2\).
For its standard projectors \(Q_b\), block expansion gives
\[
 \sum_b\norm{[\Lambda,Q_b]}_{2,N}^2
 \ge\frac{4(d-1)}{d^2}\epsred.
\]
The projection \(\Lambda\) commutes with the actual measurements.
Using \eqref{eq:physicalrounding} at \(t=1\) gives
\begin{equation}\label{eq:deltabound}
 \epsred\le\frac{dC_d^{\rm st}}{d-1}\mathfrak F_{\rm comp}.
\end{equation}
Moreover,
\[
 0\le\bar\Lambda\le I,\qquad
 \Tr\bar\Lambda=n/d,\qquad
 \epsred=\frac1{4n}\Tr(\bar\Lambda-\bar\Lambda^2).
\]
A positive contraction of trace \(k+f\), \(0\le f<1\), has
\(\Tr(\bar\Lambda-\bar\Lambda^2)\ge f(1-f)\):
moving eigenvalue mass to endpoints maximizes
\(\Tr\bar\Lambda^2\) and leaves at most one fractional eigenvalue.
Taking \(f=s/d\) and applying \eqref{eq:deltabound} and
\eqref{eq:compdefect} proves \eqref{eq:arithmeticgap}.

If \(n=dk\), choose the rank-\(k\) projection \(\bar\Lambda'\) onto
the largest eigenvalues of \(\bar\Lambda\), and put
\(\Lambda'=I_d\otimes\bar\Lambda'\).
The rank-repair estimate gives
\(\norm{\Lambda-\Lambda'}_{2,N}^2\le2\epsred\).
A closest isometry
\(\mathcal J:\C^n\to\operatorname{ran}\Lambda'\) satisfies
\[
 n^{-1}\norm{\mathcal J-\mathcal Z}_F^2\le8d\epsred.
\]
The standard projectors preserve this range. Writing
\(\RootP^t_a=\RootP^t_a(V)\), their intertwining error is
\[
 Q^t_a\mathcal J-\mathcal JP^t_a
 =(Q^t_a-\RootP^t_a)\mathcal Z
 +Q^t_a(\mathcal J-\mathcal Z)
 -(\mathcal J-\mathcal Z)P^t_a.
\]
Its total square-root norm is at most
\[
 \sqrt{224C_d^{\rm st}\mathfrak F_{\rm comp}}
 +8\sqrt{d\epsred}.
\]
Equations \eqref{eq:deltabound} and \eqref{eq:compdefect} prove
\eqref{eq:robust}. The direct bound is \(8\), two per PVM.
Finally,
\[
 C_d^{\rm st}\le2048d^2,\qquad
 \left(\sqrt{224}+\frac{8d}{\sqrt{d-1}}\right)^2<512d,
\]
giving \(K_d\le786432d^6\).

For the upper bound in \eqref{eq:gaporder}, use
\(\lfloor n/d\rfloor\) standard blocks and \(s\) one-dimensional
deterministic blocks of value \(2\).
The maximum exists by compactness of the finite-dimensional PVM space.
\end{proof}

For each fixed nonzero residue class modulo \(d\),
\eqref{eq:gaporder} gives
\(\IF-\mathcal I_{d,n}=\Theta(1/n)\).

\begin{corollary}[Rank forcing]\label{cor:ranks}
For \(d\mid n\), a projective strategy on \(\Phi_n\) satisfying
\(\Delta<2/(nK_d)\) has every outcome rank equal to \(n/d\).
\end{corollary}

\begin{proof}
For projections, \(\Tr(P-Q)^2\ge|\rank P-\rank Q|\).
Any nonconstant rank pattern has total discrepancy at least two.
Theorem~\ref{thm:robust} gives the result.
\end{proof}

\section{Noise robustness and statistical strength}\label{sec:operational}

\subsection{A local component of the standard behaviour}

Geometric summation of \eqref{eq:standard} gives
\begin{equation}\label{eq:Qd}
 Q_d(a,b|x,y)=\frac{f_d(r_{xy})}{d},\qquad
 f_d(r)=\frac1{2d^2\sin^2(\pi(r+1/4)/d)}.
\end{equation}
Construct \(L_d\) by choosing an edge uniformly, putting residual
\(d-1\) on that edge and zero on the other three, and choosing \(a_0\)
uniformly. The remaining outcomes are
\[
 b_0=a_0-r_{00},\qquad
 a_1=b_0-r_{10},\qquad
 b_1=a_1-r_{11}\pmod d.
\]
The fourth residual is consistent because the sum is \(d-1\).
This gives \(4d\) deterministic strategies of weight \(1/(4d)\).
Every edge has residual probabilities \(3/4,1/4\) at \(0,d-1\),
and all marginals are uniform. Put
\[
 a_d=4f_d(d-1)=\frac2{d^2\sin^2(3\pi/(4d))}.
\]

\begin{lemma}[Uniform-marginal local component]\label{lem:standardlocal}
For every \(d\ge2\),
\begin{equation}\label{eq:localcomponentQ}
 Q_d\ge a_dL_d,\qquad
 \frac{16}{45}<\frac{32}{9\pi^2}\le a_d\le\frac89.
\end{equation}
The coefficient \(a_d\) is maximal for this particular \(L_d\).
Furthermore,
\begin{equation}\label{eq:Qnoisebound}
 v_{\rm crit}(Q_d)\ge\frac1{2-a_d}>\frac35.
\end{equation}
\end{lemma}

\begin{proof}
With \(x=\pi/(4d)\le\pi/8\),
\[
 \frac{f_d(0)}{f_d(d-1)}
 =\left(\frac{\sin3x}{\sin x}\right)^2
 \ge(1+\sqrt2)^2>3.
\]
This gives entrywise domination. The remaining bounds in
\eqref{eq:localcomponentQ} use
\(2t/\pi\le\sin t\le t\) and \(\pi^2<10\).
The residual-\((d-1)\) entries make \(a_d\) tight for \(L_d\).

For any no-signalling behaviour \(q_{\rm rem}\) with uniform marginals,
\begin{equation}\label{eq:halfnoise}
 (q_{\rm rem}+u_d)/2\in\Loc_d.
\end{equation}
A local model guesses Alice's input uniformly, samples her guessed
output uniformly, and samples both Bob outputs conditionally according
to \(q_{\rm rem}\). Alice's other output is independent and uniform.
The guessed context gives \(q_{\rm rem}\), and the other gives \(u_d\).
Apply this to the normalized remainder of
\eqref{eq:localcomponentQ} to obtain \eqref{eq:Qnoisebound}.
In particular,
\(Q_d=\frac13L_d+\frac23q_{\rm rem}\), with a uniform-marginal
no-signalling remainder, gives a local model at visibility \(3/5\).
\end{proof}

\subsection{Tensor measurements and noise separation}

Let \(d=2^m\), identify outcomes with bit strings, and use
\(\Phi_2^{\otimes m}=\Phi_{2^m}\).
On each tensor factor measure
\[
 A_0=Z,\qquad A_1=X,\qquad
 B_0=(Z+X)/\sqrt2,\qquad B_1=(Z-X)/\sqrt2.
\]
Each party uses the same input on all factors and retains every bit.
These are two complete rank-one PVMs per party. Put
\[
 K_{xy}(a,b)=\#\{j:a_j\oplus b_j\ne xy\},\qquad
 r=\frac{2-\sqrt2}{4}.
\]
Equivalently, \(K_{xy}\) is the Hamming weight of
\(a\oplus b\oplus(xy)\ind_m\), where \(\ind_m\) is the all-ones string.
The behaviour is
\begin{equation}\label{eq:tensorbehaviour}
 T_m(a,b|x,y)=2^{-m}(1-r)^{m-K_{xy}}r^{K_{xy}}.
\end{equation}
For every deterministic local assignment, parity at each coordinate
gives
\begin{equation}\label{eq:classicalerrors}
 \sum_{x,y}K_{xy}\ge m.
\end{equation}
This applies to arbitrary dependence among the local output bits.

\begin{noiseagain}
For \(d\)-outcome projective measurements on \(\Phi_n\) with
\(I_d(q)>2\), one has
\(v_{\rm CGLMP}(q)=2/I_d(q)\ge2/\IF\), with equality precisely for
standard amplifications.
For every integer \(m\ge20000\),
\[
 v_{\rm crit}(T_m)\le\frac{17}{29}
 <\frac35<v_{\rm crit}(Q_{2^m}).
\]
\end{noiseagain}

\begin{proof}[Proof of Theorem~\ref{thm:noise}]
Since \(I_d(u_d)=0\),
\[
 I_d(vq+(1-v)u_d)=vI_d(q).
\]
Theorem~\ref{thm:value} therefore proves \eqref{eq:witnessnoise},
including the equality characterization for every ambient dimension.

For \(m\ge20000\), define the Bell functional
\[
 G_m(p)=\frac14\sum_{x,y}
 \left[\Pr_p(K_{xy}\le4m/25)-\Pr_p(K_{xy}\ge13m/25)\right].
\]
Its coefficients are \(0,\pm1/4\), and its local bound is \(1/2\).
Four low counts contradict \eqref{eq:classicalerrors}; three low counts
force the remaining count to be high.

Under \(T_m\), \(K\sim\Bin(m,r)\), where \(r<3/20\) and the variance
is \(m/8\). Each unwanted tail is bounded by
\[
 \frac{m/8}{(m/100)^2}=\frac{1250}{m}\le\frac1{16}.
\]
Thus \(G_m(T_m)\ge7/8\).
Under \(u_d\), \(K\sim\Bin(m,1/2)\), and the upper tail at \(13m/25\)
is at most
\[
 \frac{m/4}{(m/50)^2}=\frac{625}{m}\le\frac1{32}.
\]
Hence \(G_m(u_d)\ge-1/32\), and
\[
 G_m(vT_m+(1-v)u_d)\ge\frac{29v-1}{32}.
\]
Locality requires \(v\le17/29\).
Together with Lemma~\ref{lem:standardlocal}, this proves
\eqref{eq:noiseseparation}.
At \(v=3/5\), the tensor witness is at least
\[
 \frac35\frac78-\frac25\frac1{32}
 =\frac{41}{80}>\frac12,
\]
whereas \(Q_{2^m}\) has the local model above.
\end{proof}

The tensor probabilities lie in \(\mathbb Q(\sqrt2)\).
Their binomial tails can be computed from
\(b_0=(1-p)^m\) and
\(b_{k+1}=b_k(m-k)p/((k+1)(1-p))\).

\subsection{The universal statistical ceiling}

\begin{theorem}[Ceiling and closed form]\label{thm:Xi}
For every no-signalling behaviour and every fixed input law,
\begin{equation}\label{eq:Xiupper}
 \mathcal S_\mu(q)\le[1-\LF(q)]\Xi(\mu)\le\Xi(\mu).
\end{equation}
If some \(\mu_e=0\), then \(\Xi(\mu)=0\).
For full-support \(\mu\), there is a unique \(\lambda>0\) satisfying
\[
 \sum_e\min\{1,\mu_e/\lambda\}=3,
\]
and
\begin{equation}\label{eq:Xiformula}
 \Xi(\mu)=\sum_{\mu_e<\lambda}
             \mu_e\log_2(\lambda/\mu_e).
\end{equation}
Moreover,
\(\Xi(\mu)\le\cKL:=\log_2(4/3)\), with equality exactly for uniform
inputs.
\end{theorem}

\begin{proof}
Delete an edge \(e\) from the four-cycle.
The remaining three consistent pair distributions extend to a joint
distribution on all four outputs by conditional gluing. Deleting
\(11\), for example, gives
\[
 q(a_0,b_0|00)\,q(a_1|b_0,10)\,q(b_1|a_0,01).
\]
Conditionals on zero-marginal events may be chosen arbitrarily.
Let \(L^{(e)}\) be the resulting local extension.
For an optimizer \(s\) in \eqref{eq:Xi}, mix these extensions with
weights \(1-s_e\). The local behaviour \(L_*\) satisfies
\((L_*)_e\ge s_eq_e\), hence
\(D_\mu(q\Vert L_*)\le\Xi(\mu)\).
Apply this comparison to the no-signalling remainder in
\(q=\alpha L+(1-\alpha)q_{\rm rem}\), then use joint convexity of
relative entropy to obtain \eqref{eq:Xiupper}.

The minimum is unchanged by replacing \(\sum_es_e=3\) with
\(\sum_es_e\le3\), since the objective decreases in each sampled
coordinate. If an input is absent, choose \(s_e=1\) on all sampled
contexts. For full support, the convex optimality conditions give
\(s_e=\min(1,\mu_e/\lambda)\), proving \eqref{eq:Xiformula}.
To evaluate it, sort the four weights, test capped sets of sizes
\(k=0,1,2\), and use
\[
 \lambda=\frac{\sum_{\rm uncapped}\mu_e}{3-k}.
\]
Choose the set for which capped weights are at least \(\lambda\)
and uncapped weights at most \(\lambda\); ties contribute zero.
Finally \(s_e=3/4\) gives \(\Xi(\mu)\le\cKL\).
For nonuniform \(\mu\), a tangent perturbation at this interior point
strictly decreases the objective.
\end{proof}

\subsection{Positive local content of quantum behaviours}

\begin{theorem}[Local-component bound]\label{thm:LF}
For arbitrary finite-dimensional states and POVMs in the
\((2,2,d)\) scenario,
\begin{equation}\label{eq:gamma}
 \LF(q)\ge\gamma_d:=d-\sqrt{d^2-1}
 =\frac1{d+\sqrt{d^2-1}}>\frac1{2d}.
\end{equation}
More generally, if Alice's two measurements have \(m_0,m_1\) outcomes,
the bound is
\[
 \frac{m_A-\sqrt{m_A^2-4}}2,\qquad m_A=m_0+m_1.
\]
\end{theorem}

This explicit dependence on the outcome counts continues the
two-input positive-content results of \cite{GMS,Joshi}.

\begin{proof}
Purify the state and dilate the finite POVMs locally to PVMs.
The dilation is the isometry
\(\xi\mapsto\sum_a\sqrt{F_a}\xi\otimes\ket a\), extended to a unitary
separately for each setting with a common initial ancilla.
We may therefore use a vector \(\psi\), Alice projections \(A_a,C_c\),
and Bob projections \(B_b^y\).

Set
\[
 \begin{gathered}
 u_{ab}^y=A_aB_b^y\psi,\qquad
 v_{cb}^y=C_cB_b^y\psi,\qquad
 p_{ab}^y=\norm{u_{ab}^y}^2,\qquad
 r_{cb}^y=\norm{v_{cb}^y}^2,\\
 z_{acb}^y=\ip{u_{ab}^y}{v_{cb}^y},\qquad
 k_{ac}=\ip{A_a\psi}{C_c\psi},\qquad
 w_{ac}=q_A(a|0)+q_A(c|1).
 \end{gathered}
\]
Then
\[
 \sum_bz_{acb}^y=k_{ac},\qquad
 \sum_{a,c}k_{ac}=1,\qquad
 \sum_{a,c}w_{ac}=m_A.
\]
Define, with \(0/0=0\),
\[
 t_{acb}^y=\frac{|z_{acb}^y|^2}
                  {\max\{p_{ab}^y,r_{cb}^y\}}.
\]
At fixed \(y,b\), the nonzero vectors
\(v_{cb}^y/\sqrt{r_{cb}^y}\) are orthonormal in \(c\).
Bessel's inequality and its counterpart in \(a\) give
\begin{equation}\label{eq:capacities}
 \sum_ct_{acb}^y\le p_{ab}^y,\qquad
 \sum_at_{acb}^y\le r_{cb}^y.
\end{equation}

Fix \(a,c\), abbreviate \(w=w_{ac}\), \(k=|k_{ac}|\), and put
\[
 h_{ac}=\frac{w-\sqrt{w^2-4k^2}}2,
\]
with value zero if \(w=0\).
Cauchy--Schwarz gives \(k\le w/2\).
For \(T_y=\sum_bt_{acb}^y\) and
\(z_y=\sum_b\min(p_{ab}^y,r_{cb}^y)\), one has \(z_y\le w/2\) and
\[
 k\le\sum_b\sqrt{p_{ab}^yr_{cb}^y}
 \le\sqrt{z_y(w-z_y)}.
\]
Thus \(z_y\ge h_{ac}\). Weighted Cauchy--Schwarz gives, for \(k>0\),
\[
 T_y\ge\frac{k^2}{w-z_y}
 \ge\frac{k^2}{w-h_{ac}}=h_{ac};
\]
the zero case follows directly.

Assign the deterministic strategy \((a,c,b_0,b_1)\) weight
\begin{equation}\label{eq:localweights}
 w^{\rm loc}(a,c,b_0,b_1)=
 h_{ac}\frac{t_{acb_0}^0}{T_0}\frac{t_{acb_1}^1}{T_1},
\end{equation}
with value zero for \(h_{ac}=0\).
Otherwise the denominators are positive.
By \eqref{eq:capacities} and \(h_{ac}\le T_y\), the weights
\eqref{eq:localweights} define a subnormalized local behaviour
entrywise below \(q\), of mass \(\alpha=\sum_{a,c}h_{ac}\).

The function \(\phi(t)=(1-\sqrt{1-4t^2})/2\) is convex and increasing
on \([0,1/2]\). Weighted Jensen, omitting zero weights, yields
\[
 \alpha=\sum_{a,c}w_{ac}\phi(|k_{ac}|/w_{ac})
 \ge m_A\phi\!\left(\frac{\sum_{a,c}|k_{ac}|}{m_A}\right)
 \ge m_A\phi(1/m_A),
\]
where \(\sum|k_{ac}|\ge1\) was used in the last step.
This proves the general bound, and \(m_A=2d\) gives \eqref{eq:gamma}.
\end{proof}

Theorems~\ref{thm:Xi} and \ref{thm:LF} imply, uniformly in the state
dimension,
\begin{equation}\label{eq:nonattainmentbound}
 \mathcal S_\mu(q)\le(1-\gamma_d)\Xi(\mu)
\end{equation}
for every finite-outcome quantum behaviour.
For full-support \(\mu\), \(\Xi(\mu)>0\), giving a strict gap.
If a context is absent, the sampled graph is a forest, and every
no-signalling behaviour has zero strength on the sampled data.

\subsection{The supremum and an explicit separation}

\begin{supremumagain}
For every fixed input law independent of the hidden variable,
\[
 \sup_{\substack{d\ge2\\
     \text{two rank-one PVMs per party on }\Phi_d}}
 \mathcal S_\mu(q)=\Xi(\mu).
\]
The same supremum holds over arbitrary finite-dimensional quantum
states and finite-outcome POVMs. For full-support \(\mu\), every
finite-outcome quantum behaviour has strength strictly below
\(\Xi(\mu)\).
\end{supremumagain}

\begin{proof}[Proof of Theorem~\ref{thm:supremum}]
Use \(T_m\) from \eqref{eq:tensorbehaviour} and the strict events
\(\mathcal A_{xy}=\{K_{xy}<m/4\}\), with common quantum probability
\[
 t_m=\Pr\{\Bin(m,r)<m/4\}.
\]
By \eqref{eq:classicalerrors}, at most three of the four events hold
on a deterministic local assignment. Thus local event probabilities
obey \(\sum_ez_e\le3\).
Data processing, retaining the context and both the event and its
complement, gives
\begin{equation*}
 \mathcal S_\mu(T_m)\ge
 \min_{\substack{0\le z_e\le1\\\sum_ez_e\le3}}
 \sum_e\mu_eD_{\rm bin}(t_m\Vert z_e)
 \ge t_m\Xi(\mu)-h_2(t_m).
\end{equation*}
The final inequality drops the nonnegative terms
\(-(1-t_m)\log_2(1-z_e)\).

A centered Bernoulli log moment-generating function has second
derivative at most \(1/4\). Integrating twice and optimizing the
Chernoff bound gives
\[
 1-t_m\le e^{-\kappa m},\qquad
 \kappa=2(1/4-r)^2=\frac{3-2\sqrt2}{8}>\frac1{50},
\]
where the last comparison follows from \((71/50)^2>2\).
Using \(h_2(\epsilon)\le\epsilon\log_2(e/\epsilon)\), we obtain
\begin{equation}\label{eq:KLrate}
 0\le\Xi(\mu)-\mathcal S_\mu(T_m)
 \le e^{-\kappa m}
 \left[\Xi(\mu)+\frac{\kappa m+1}{\ln2}\right].
\end{equation}
This tends to zero. Theorem~\ref{thm:Xi} supplies the universal
upper bound, and \eqref{eq:nonattainmentbound} supplies the
finite-outcome quantum gap, proving \eqref{eq:supremum}.
\end{proof}

The binary PR behaviour attains \(\Xi(\mu)\) in \(\NS_2\):
its four parity tests succeed with certainty, while local success
probabilities satisfy \(\sum_ez_e\le3\).
For uniform inputs this gives the value computed in \cite{Grudka}.

From \eqref{eq:localcomponentQ} and \eqref{eq:Xiupper},
\begin{equation}\label{eq:QKL}
 \mathcal S_\mu(Q_d)\le(1-a_d)\Xi(\mu),\qquad
 \sup_\mu\mathcal S_\mu(Q_d)<\frac{29}{45}\log_2\frac43.
\end{equation}
For \(m\ge200\), \(\kappa m>4\).
The function \(e^{-x}[\cKL+(x+1)/\ln2]\) decreases for \(x\ge0\), and
\[
 \ln(4/3)=2\artanh(1/7)>\frac27,\qquad
 e^4>\sum_{j=0}^9\frac{4^j}{j!}>54.
\]
Equation \eqref{eq:KLrate} therefore gives
\[
 \frac{\cKL-\mathcal S_{\rm unif}(T_m)}{\cKL}
 <\frac1{54}\left(1+\frac5{2/7}\right)=\frac{37}{108}.
\]
Since \(71/108=29/45+7/540\), combining this with \eqref{eq:QKL}
proves, for every \(m\ge200\),
\begin{equation}\label{eq:KLseparation}
 \mathcal S_{\rm unif}(T_m)>
 \frac{71}{108}\log_2\frac43>
 \frac{29}{45}\log_2\frac43>
 \sup_\mu\mathcal S_\mu(Q_{2^m}).
\end{equation}

\subsection{Statistical strength near CGLMP optimality}

For behaviours in the same finite scenario,
\begin{equation}\label{eq:LFLipschitz}
 |\LF(p)-\LF(q)|\le\sum_{x,y}\TV(p_{xy},q_{xy}).
\end{equation}
To prove this, write a sublocal component \(F\le p\) with deterministic
weights \(w_v\). For each event \(e\), set
\(t_e=(F_e-q_e)_+/F_e\), with value zero if \(F_e=0\), and replace
\[
 w_v\longmapsto
 w_v\prod_{\text{contexts }c}(1-t_{c,v(c)}).
\]
The resulting sublocal component lies below \(q\).
The union bound controls its lost mass by
\[
 \sum_e(F_e-q_e)_+\le\sum_c\TV(p_c,q_c).
\]
Interchanging \(p,q\) proves \eqref{eq:LFLipschitz}.

Let \(P\) be \(d\)-outcome PVMs on \(\Phi_{dk}\), and choose a standard
reference \(Q\) realizing \eqref{eq:distance}.
Write \(\epsmeas=\epsmeas(P,Q)\).
For each context, the identity
\(\sqrt{\tau(PR)}=\norm{PR}_2\), completeness, and the triangle
inequality give
\[
 H(q_{xy},Q_{d,xy})\le\sqrt{e_{2x}}+\sqrt{e_{1+2y}},
 \qquad e_t=\sum_a\norm{P^t_a-Q^t_a}_2^2.
\]
Since \(\TV\le H\),
\(\sum_{x,y}\TV(q_{xy},Q_{d,xy})\le4\sqrt{\epsmeas}\).
Theorem~\ref{thm:robust} and \eqref{eq:LFLipschitz}, with
\(\widehat K_d=786432d^6\), imply
\[
 \LF(q)\ge
 \bigl(a_d-4\sqrt{\widehat K_d\Delta}\bigr)_+.
\]
Together with \eqref{eq:gamma},
\begin{equation}\label{eq:KLtradeoff}
 \mathcal S_\mu(q)\le
 \left[
 1-\max\left\{\gamma_d,
       \bigl(a_d-4\sqrt{\widehat K_d\Delta}\bigr)_+\right\}
 \right]\Xi(\mu).
\end{equation}
The uniform-noise construction \eqref{eq:halfnoise} uses the
uniform-marginal component \(L_d\); the local-fraction estimates use
entrywise domination for general marginals.

\section{Discussion}\label{sec:discussion}

For projective measurements on maximally entangled states, maximal
CGLMP violation determines the Fourier-phase measurements and their
multiplicity.
The quantitative bounds also detect the arithmetic obstruction to
equality and give an optimal deficit of order \(1/n\) along fixed
nonzero residue classes.
The operational comparisons distinguish optimization of the CGLMP
witness from global noise resistance and statistical strength.
The function \(\Xi(\mu)\) determines the dimension-unrestricted
statistical ceiling, while positive local content gives a finite-outcome
gap and \eqref{eq:KLtradeoff} relates that gap to projective CGLMP
rigidity.

Several optimization problems remain open. For \(d\)-outcome POVMs
on a fixed \(\Phi_n\), determine the maximum CGLMP value and classify
the maximizing measurements. For optimization over states, determine
the sharp value and all equality cases at general \(d\), extending
the low-dimensional analytic results \cite{IR}.
When \(n\bmod d>0\), determine the projective optimum and the
structure of its maximizers beyond the order estimate
\eqref{eq:gaporder}.

For a prescribed finite outcome count and input law, determine the
statistical-strength optimum and its optimizing state and
measurements. A related question is to determine the sharp
outcome-dependent local-fraction bound, including its dependence
on unequal outcome counts. These problems would sharpen both the
finite-dimensional ceiling gap and the relation between Bell-value
optimality and statistical evidence.

\section*{Code availability}

Scripts executing the rational interval and Bernstein-coefficient
calculations, and evaluating the integer facet minor, are available
from the authors. They use the data and procedures described in
Appendices~\ref{app:facet}--\ref{app:continuous}.

\appendix

\section{The Bancal--Gisin--Pironio four-outcome facet}\label{app:facet}

Bancal, Gisin, and Pironio \cite[Appendix B]{BGP} found the inequality
\(S^2_{(2,2,4)}\) as a genuinely four-outcome facet. We give its
integer rank certificate, a rational no-signalling point, and the
outcome-fusion argument.

Use each party's six independent marginals, for example
\[
 p_A=(p_A(0|0),p_A(1|0),p_A(2|0),
      p_A(0|1),p_A(1|1),p_A(2|1)),
\]
and the \(36\) joint probabilities with \(a,b<3\).
The affine dimension is \(48\). The inequality is
\begin{equation}\label{eq:facet}
 F(p)=\ell\cdot p_A+\ell\cdot p_B+
 \sum_{x,y=0}^1\sum_{a,b=0}^2
 H_{3x+a,\,3y+b}p(a,b|x,y)\le0,
\end{equation}
where
\[
 \ell=(-1,-1,0,-1,0,0),\qquad
 H=\begin{pmatrix}
 0&1&0&1&0&1\\
 1&0&0&1&1&0\\
 0&0&-1&1&0&0\\
 1&1&1&1&0&0\\
 0&1&0&0&0&-1\\
 1&0&0&0&-1&0
 \end{pmatrix}.
\]

\begin{proposition}\label{prop:bgp}
Inequality \eqref{eq:facet} defines a facet of \(\Loc_4\) meeting
the relative interior of \(\NS_4\). Its class is distinct from
CGLMP\(_4\) and from every outcome-fusion lifting of a lower-outcome
inequality. There is a rational no-signalling behaviour \(p^+\)
with \(F(p^+)=1/144\).
Four is the smallest outcome count for such a facet.
\end{proposition}

\begin{proof}
For \(s=(a_0,a_1,b_0,b_1)\), define
\[
 A(s)=(\ind_{a_0=0},\ind_{a_0=1},\ind_{a_0=2},
       \ind_{a_1=0},\ind_{a_1=1},\ind_{a_1=2}),
\]
and define \(B(s)\) analogously. Then
\[
 F(s)=\ell\cdot A(s)+\ell\cdot B(s)+A(s)^THB(s).
\]
Evaluation on the \(4^4\) strategies gives
Table~\ref{tab:facetdata}, proving local validity.
Let \(\mathcal S=\{s:F(s)=0\}\), and augment the coordinates:
\[
 z(s)=\bigl(1,A(s),B(s),(A_i(s)B_j(s))_{i,j=0}^5\bigr),
 \qquad \iota(s)=64a_0+16a_1+4b_0+b_1,
\]
with the products in row-major order.
Select the rows with the ordered identifiers in
Table~\ref{tab:minor}. Numbering columns from zero and deleting
column \(47\) gives a \(48\times48\) matrix \(Z_{\rm min}\) with
\begin{equation}\label{eq:minor}
 \det Z_{\rm min}=-1.
\end{equation}
Thus the saturation face has affine rank \(47\).

The average
\[
 \bar p=\frac1{72}\sum_{s\in\mathcal S}q_s
\]
has minimum entry \(1/72\), obtained by counting occurrences of
each event among the saturators.
Since probability nonnegativity defines \(\NS_4\) in its affine
hull, \(\bar p\) is a relative-interior point.
Modify only context \(00\), adding \(1/144\) at \((0,1),(3,3)\)
and subtracting \(1/144\) at \((0,3),(3,1)\).
The resulting \(p^+\) preserves all row and column sums, remains
nonnegative, and satisfies \(F(p^+)=1/144\).

A deterministic CGLMP\(_d\) saturator is specified by an anchor
outcome and four nonnegative residuals summing to \(d-1\).
Its saturation count is \(d\binom{d+2}{3}\).
At \(d=4\), the counts \(80\) and \(72\) distinguish the two classes.

Two outcomes are twins when replacing one by the other preserves
the functional for every assignment of the other three outcomes.
Every nontrivial fusion lifting has twins.
Here the four extended coefficient rows, including the zero row
for output \(3\), together with their marginal coefficients, are
distinct within each setting. The analogous column profiles are
also distinct. Taking the other party's unused outputs to be \(3\)
isolates any differing coefficient. This distinguishes every
outcome pair and separates the facet from all nontrivial fusion
liftings.

The complete two-outcome classification is due to Fine
\cite{Fine}; the three-outcome list uses the classifications and
facet results of \cite{CG,Masanes}.
Their counts in Table~\ref{tab:smallfacets} show that \(d=4\) is the
first outcome count with the asserted properties.
\end{proof}

\begin{table}[ht]
\centering
\caption{Deterministic values of \eqref{eq:facet}.}
\label{tab:facetdata}
\begin{tabular}{@{}lrrr@{}}
\toprule
\(F(s)\) & \(0\) & \(-1\) & \(-2\)\\
\midrule
Number of strategies & \(72\) & \(144\) & \(40\)\\
\bottomrule
\end{tabular}
\end{table}

\begin{table}[ht]
\centering
\small
\caption{Ordered saturator identifiers for the minor
\eqref{eq:minor}.}
\label{tab:minor}
\begin{tabular}{@{}rrrrrrrrrrrr@{}}
\toprule
4&6&8&20&21&22&23&34&38&42&46&54\\
58&62&64&65&72&81&85&89&93&96&97&98\\
99&113&121&125&128&132&136&140&149&151&156&162\\
163&172&188&191&200&201&202&203&213&215&221&226\\
\bottomrule
\end{tabular}
\end{table}

\begin{table}[ht]
\centering
\caption{Complete facet counts for two and three outcomes.}
\label{tab:smallfacets}
\begin{tabular}{@{}crrr@{}}
\toprule
\(d\) & Positivity & Fused CHSH & CGLMP\(_3\)\\
\midrule
2 & 16 & 8 & 0\\
3 & 36 & 648 & 432\\
\bottomrule
\end{tabular}
\end{table}

For general \(d\), the independent coordinates are each party's
\(2(d-1)\) marginals and the \(4(d-1)^2\) joint probabilities with
\(a,b<d-1\). For example,
\[
 p(a,d-1|x,y)=p_A(a|x)-\sum_{b<d-1}p(a,b|x,y),
\]
with the analogous Bob formula and normalization determining the
remaining entries. Mixed differences of deterministic vertices show
full affine span.

The small-outcome lists can also be obtained by the rational
double-description method \cite{FukudaProdon}.
With \(L=4d(d-1)+1\), start from the columns of \(Z^{-1}\), where
\(Z\) is a greedily selected set of \(L\) independent augmented rows.
Insert the other rows lexicographically. For a new row \(z\),
retain rays with \(zr\ge0\), and join a positive and negative ray by
\((zr_+)r_--(zr_-)r_+\) when their common active rows have rank
\(L-2\). This rank condition identifies a two-dimensional face.
Primitive-integer normalization removes duplicate rays.
Comparing final value vectors, up to positive proportionality,
with positivity, relabelled CGLMP functionals, and CGLMP\(_2\)
composed with output surjections gives Table~\ref{tab:smallfacets}.
For four outcomes, the larger list of Jesus and
Zambrini Cruzeiro contains \(34\) classes and \(11\,665\,992\)
facets \cite{JC}; adjacency sampling provides a complementary
enumeration method \cite{Staufenbiel}.

\section{The circulation sum of squares}\label{app:circulation}

Fix \(i<j<k\), and write
\[
 \begin{array}{lll}
 u_\ell=E_{i\ell}E_{j,\ell+2},&
 v_\ell=E_{j\ell}E_{k,\ell+2},&
 z_\ell=E_{i\ell}E_{k,\ell+2},\\
 a_\ell=E_{i\ell}E_{j,\ell+1},&
 b_\ell=E_{j\ell}E_{k,\ell+1},&
 c_\ell=E_{i\ell}E_{k,\ell+1}.
 \end{array}
\]
Let \(P\) be the sum of squared norms of \(a_\ell,b_\ell,c_\ell\),
and \(O\) the corresponding sum for \(u_\ell,v_\ell,z_\ell\).
Then \(P+O=p_{ij}+p_{jk}+p_{ik}\).
Use the triples in Table~\ref{tab:circulation} for each
\(\ell\in\Z_4\).

\begin{table}[ht]
\centering
\caption{Terms in the circulation identity.}
\label{tab:circulation}
\begin{tabular}{@{}ccc@{}}
\toprule
\(X_{\ell\nu}\) & \(Y_{\ell\nu}\) & \(\epsilon_\nu\)\\
\midrule
\(z_\ell\) & \(a_\ell\) & \(1\)\\
\(z_\ell^\dagger\) & \(b_{\ell+1}^\dagger\) & \(1\)\\
\(v_{\ell+1}\) & \(a_\ell^\dagger\) & \(1\)\\
\(u_\ell\) & \(c_\ell\) & \(-1\)\\
\(u_\ell^\dagger\) & \(b_{\ell+2}\) & \(1\)\\
\(v_{\ell+3}^\dagger\) & \(c_\ell^\dagger\) & \(-1\)\\
\bottomrule
\end{tabular}
\end{table}

Projection completeness gives
\[
 \begin{split}
 C_{ijk}={}&-2\sum_\ell\norm{z_\ell}_2^2
 +2\Re\sum_\ell
       \bigl[\ip{u_\ell}{z_\ell}
          +\ip{v_\ell^\dagger}{z_\ell^\dagger}\bigr]\\
 &+2\Re\sum_{\ell,\nu}
       \epsilon_\nu\ip{X_{\ell\nu}}{Y_{\ell\nu}}.
 \end{split}
\]
For \(\sigma=\pm1\),
\begin{equation}\label{eq:circulationSOS}
 \begin{split}
 2P+5O-\sigma C_{ijk}
 ={}&\sum_{\ell,\nu}
       \norm{Y_{\ell\nu}-\sigma\epsilon_\nu X_{\ell\nu}}_2^2\\
 &+3\sum_\ell\norm{u_\ell-\tfrac\sigma3z_\ell}_2^2
 +3\sum_\ell
       \norm{v_\ell^\dagger-\tfrac\sigma3z_\ell^\dagger}_2^2\\
 &+\left(\frac73+2\sigma\right)
       \sum_\ell\norm{z_\ell}_2^2\ge0.
 \end{split}
\end{equation}
Thus \(|C_{ijk}|\le2P+5O\le5(P+O)\).

Both identities follow by expanding \(E_{j3}=I-E_{j0}-E_{j1}-E_{j2}\),
reducing adjacent same-site products by orthogonality, and comparing
coefficients of cyclic words. Choose the lexicographically least
cyclic rotation of each word, and expand each real part into the
half-sum of a word and its adjoint, keeping those coefficients
separate.

\section{Scalar coefficients and the alternative estimate}
\label{app:scalar}

\subsection{Coefficient and interval data}

The coefficients of
\(r_{\rm reg}(y)=\csc(\pi y/2)-2/(\pi y)\) are
\begin{equation}\label{eq:rn}
 r_n=
 \frac{2(2^{2n-1}-1)|B_{2n}|}{(2n)!}
 \left(\frac\pi2\right)^{2n-1},
 \qquad 0<r_n<3\,4^{-n},
 \qquad
 r_{\rm reg}(y)=\sum_{n\ge1}r_ny^{2n-1}.
\end{equation}
For \(y=1-h/d\), let
\(m_k=\frac d{d+1}(d-h+1)_k/(d+2)_k\).
The error in \eqref{eq:es} is
\begin{equation}\label{eq:finiteeseries}
 e_{d,h}=\sum_{n\ge1}r_n
 \left[
 y^k-\left(1+\frac{(k+1)(k+2)}{2d}\right)m_k
 +\frac{k(k+1)}{2d}m_{k-1}
 \right],\qquad k=2n-1.
\end{equation}
The bracket is bounded in absolute value by
\(2+4n^2/(d+1)\). Truncation after \(32\) terms has error at most
\begin{equation}\label{eq:finitetail}
 T_d=3\sum_{n=33}^\infty4^{-n}
           \left(2+\frac{4n^2}{d+1}\right).
\end{equation}
For \(q=1/4\), the tail is evaluated using
\[
 \sum_{n\ge J}q^n=\frac{q^J}{1-q},\qquad
 \sum_{n\ge J}n^2q^n
 =q^J\left[
 \frac{J^2}{1-q}+\frac{2Jq}{(1-q)^2}
 +\frac{q(1+q)}{(1-q)^3}
 \right].
\]

For the coefficient enclosures, put
\[
 A(k,N)=\sum_{j=0}^{N-1}\frac{(-1)^j}{(2j+1)k^{2j+1}}.
\]
Machin's identity gives the interval
\[
 \pi\in16[A(5,128),A(5,129)]
       -4[A(239,128),A(239,129)].
\]
Generate \(B_0=1\) and
\(B_m=-(m+1)^{-1}\sum_{j=0}^{m-1}\binom{m+1}{j}B_j\)
through \(m=64\), and substitute the interval for \(\pi\) into
\eqref{eq:rn}.
For \eqref{eq:finiteeseries}, use
\(m_0=d/(d+1)\) and
\(m_k=m_{k-1}(d-h+k)/(d+k+1)\).
Endpoint inclusion in these rational operations, followed by
\([-T_d,T_d]\), encloses each \(e_{d,h}\).
For interval triangle loads in \eqref{eq:load}, use
\(\amag([l,u])=\max(|l|,|u|)\).
The parameters and resulting lower bound are collected in
Table~\ref{tab:intervals}.

\begin{table}[ht]
\centering
\small
\caption{Data for the finite interval calculation. The margin uses
lower endpoints for \(e_{d,h}\) and upper bounds for
\(\ell_{d,h}\).}
\label{tab:intervals}
\begin{tabular}{@{}ll@{}}
\toprule
Quantity & Data\\
\midrule
Outcome counts & \(2\le d\le65\)\\
Separations & \(1\le h<d\)\\
Cosecant terms & \(1\le n\le32\)\\
Moment indices & \(0\le k\le63\)\\
Bernoulli indices & \(0\le m\le64\)\\
Arctangent partial sums & \(128,\ 129\) terms\\
Tail start in \eqref{eq:finitetail} & \(J=33\)\\
Lower bound for the normalized margin & \(1935/10000>1/6\)\\
\bottomrule
\end{tabular}
\end{table}

Consequently, for every \(2\le d\le65\),
\begin{equation}\label{eq:finitecondition}
 (d+1)\left[
 \frac1{3(d+1)}+\min_he_{d,h}
 -\frac5{2d}\max_h\ell_{d,h}
 \right]>\frac16.
\end{equation}
This is the finite-range input to \eqref{eq:piecewisecertificate}.

\subsection{The beta--Stein estimate}

Put \(\varphi(y)=y(1-y)r_{\rm reg}(y)\) and
\(Z\sim\Beta(dy,d(1-y))\). For \(y=1-h/d\),
\begin{equation}\label{eq:betaerror}
 e_{d,h}=
 \frac{\varphi(y)-\mathbb E\varphi(Z)
       +\tfrac12\mathbb E[(Z-y)\varphi'(Z)]}{y(1-y)}.
\end{equation}
Indeed, the pole cancels because
\(B_d((1-x)^{-1})(h)=d/(d-h)\).
Size bias replaces \(B_dg\) by
\(\mathbb E[Z(1-Z)g(Z)]/[y(1-y)]\), and
\[
 \mathbb E[(Z-y)g(Z)]
 =d^{-1}\mathbb E[Z(1-Z)g'(Z)].
\]
These identities give \eqref{eq:betaerror}.

For \(k=2n-1\), the third and fourth derivatives of
\(y^{k+1}-y^{k+2}\) are bounded by \(3k(k+1)\) and
\(4k(k-1)(k+1)\). Summing the positive series \eqref{eq:rn} and using
\[
 r_{\rm reg}''(1)+2r_{\rm reg}'(1)=\pi^2/4,\qquad
 r_{\rm reg}'''(1)+3r_{\rm reg}''(1)=3\pi^2/4
\]
gives
\(\norm{\varphi'''}_\infty\le3\pi^2/4\) and
\(\norm{\varphi^{(4)}}_\infty\le3\pi^2\).
The trapezoidal remainder is
\[
 \begin{split}
 \varphi(y)-\varphi(y+\xi)
 +\frac\xi2[\varphi'(y+\xi)+\varphi'(y)]
 &=\frac{\xi^3}{12}\varphi'''(y)+\epsilon_\xi,\\
 |\epsilon_\xi|&\le
 \frac{\norm{\varphi^{(4)}}_\infty}{24}|\xi|^4.
 \end{split}
\]
For \(\xi=Z-y\),
\[
 \mathbb E\xi^3=
 \frac{2y(1-y)(1-2y)}{(d+1)(d+2)},\qquad
 \mathbb E\xi^4=
 \frac{3y(1-y)[2(1-2y)^2+y(1-y)(d+2)]}
 {(d+1)(d+2)(d+3)}.
\]
Thus, with \(q=|1-2y|\),
\[
 |e_{d,h}|\le
 \frac{\pi^2}{8(d+1)(d+2)}
 \left[q+\frac{3(d+2)}{4(d+3)}
 -\frac{3(d-6)}{4(d+3)}q^2\right].
\]
For \(d\ge66\), the bracket is at most
\[
 \frac{3(d+2)}{4(d+3)}+\frac{d+3}{3(d-6)}<\frac98,
\]
the difference from \(9/8\) being
\((d^2-57d-342)/(24(d-6)(d+3))>0\).
Using \(\pi^2<10\), we obtain
\begin{equation}\label{eq:largeM}
 b_d\le\frac{45}{32(d+1)(d+2)}.
\end{equation}

To derive \eqref{eq:piecewisecertificate}, combine
\eqref{eq:largeM}, \eqref{eq:coarseerror}, and
\(L_s\ge2\Fdef/[3(d+1)]\) with \eqref{eq:matching}.
For \(2\le d\le65\), instead use \eqref{eq:Hodge},
\eqref{eq:circulation}, and \eqref{eq:finitecondition}:
\[
 L_{e+s}\ge
 2\left[
 \frac1{3(d+1)}+\min_he_{d,h}
 -\frac5{2d}\max_h\ell_{d,h}
 \right]\Fdef
 \ge\frac{\Fdef}{3(d+1)}.
\]

\section{Polynomial positivity on the continuous domain}
\label{app:continuous}

For \(\theta\ge0\), define
\[
 Q_{k,\theta}(y)=
 \prod_{j=1}^k\frac{y+j\theta}{1+(j+1)\theta},
 \qquad Q_{0,\theta}=1,
\]
and
\[
 \mathcal E_{k,\theta}(y)=
 \frac{(1+\theta)y^k-Q_{k,\theta}(y)}{\theta}
 -\frac{(k+1)(k+2)}2Q_{k,\theta}(y)
 +\frac{k(k+1)}2Q_{k-1,\theta}(y),
\]
with the removable value at zero filled continuously.
Put \(\mathcal E_\theta(y)=\sum_{n\ge1}r_n\mathcal E_{2n-1,\theta}(y)\).
Formula \eqref{eq:finiteeseries} gives
\begin{equation}\label{eq:Enormalization}
 \mathcal E_{1/d}(1-h/d)=(d+1)e_{d,h}.
\end{equation}
The required scalar bound is
\begin{equation}\label{eq:continuouslemma}
 \begin{gathered}
 \frac13+\mathcal E_\theta(y)
 -\frac52|\Omega_\theta(x_1,x_2,x_3)|\ge\frac16,\\
 \Omega_\theta=(1-2\theta)
 \sum_{j=1}^3[\mathcal E_\theta(1-x_j)-\mathcal E_\theta(x_j)],
 \end{gathered}
\end{equation}
for \(0<\theta\le1/3\), \(\theta\le y\le1-\theta\),
\(x_j\ge\theta\), and \(\sum_jx_j=1\).

\subsection{Truncation data}

Let \(Q=10^{12}\), and take the integers in
Table~\ref{tab:cosecant}.
The intervals from Appendix~\ref{app:scalar} give
\(m_n/Q<r_n<(m_n+1)/Q\).
The partial-fraction expansion of cosecant yields
\[
 r_n=\frac4\pi\,4^{-n}\eta(2n)<\frac43\,4^{-n},
 \qquad
 \eta(2n)=\sum_{j\ge1}(-1)^{j-1}j^{-2n}\in(0,1).
\]
Each factor
\(q_j=(y+j\theta)/(1+(j+1)\theta)\) belongs to \([0,1]\) for
\(0\le y\le1\), with \(|q_j-y|\le j\theta\).
Hence \(|\mathcal E_{k,\theta}(y)|\le1+(k+1)^2\).
Replacing \(r_n\) by \(m_n/Q\) for \(n\le10\) and omitting the
remaining terms has uniform error at most
\begin{equation}\label{eq:uniformtail}
 \epssc=
 \frac{1550}{10^{12}}
 +\frac43\sum_{n=11}^{\infty}4^{-n}(1+4n^2)
 =\frac{45400712033}{207360000000000}
 <\varepsilon_0:=\frac9{40000}.
\end{equation}

\begin{table}[ht]
\centering
\caption{Integers specifying \(m_n/Q<r_n<(m_n+1)/Q\), with
\(Q=10^{12}\).}
\label{tab:cosecant}
\begin{tabular}{@{}rrrr@{}}
\toprule
\(n\) & \(m_n\) & \(n\) & \(m_n\)\\
\midrule
1 & 261799387799 & 6 & 310774174\\
2 & 75362478042  & 7 & 77707647\\
3 & 19606915981  & 8 & 19427797\\
4 & 4954856459   & 9 & 4857004\\
5 & 1242203718   & 10 & 1214254\\
\bottomrule
\end{tabular}
\end{table}

\subsection{The polynomials and their Bernstein bounds}

Let
\[
 D_*(\theta)=\prod_{j=2}^{20}(1+j\theta),\qquad
 D_k=\prod_{j=2}^{k+1}(1+j\theta),\qquad
 P_k=\prod_{j=1}^k(y+j\theta),
\]
with empty products one, and set
\(a_k=(k+1)(k+2)/2\), \(b_k=k(k+1)/2\).
Define
\begin{equation*}
 \begin{split}
 N(\theta,y)
 &=\sum_{\substack{1\le n\le10\\k=2n-1}}
 m_n\frac{D_*}{D_k}\frac{\mathcal B_k(\theta,y)}{\theta},\\
 \mathcal B_k(\theta,y)
 &=(1+\theta)y^kD_k-(1+a_k\theta)P_k
   +b_k\theta P_{k-1}(1+(k+1)\theta).
 \end{split}
\end{equation*}
Every \(\mathcal B_k\) is divisible by \(\theta^2\).
Thus \(N\) is an integer polynomial of bidegree at most \((19,19)\),
and the truncated function is \(N/(QD_*)\).

Parameterize the domains by
\[
 Y=\theta+(1-2\theta)u,\qquad
 X_j=\theta+(1-3\theta)\alpha_j,\qquad
 0\le u\le1,\quad\alpha_j\ge0,\quad\sum_j\alpha_j=1,
\]
and put
\[
 \mathcal N_\triangle=(1-2\theta)
 \sum_{j=1}^3[N(\theta,1-X_j)-N(\theta,X_j)].
\]
The polynomial inequalities are
\begin{equation}\label{eq:rectangle}
 100N(\theta,Y)+3QD_*(\theta)>0,
 \qquad 0\le\theta\le1/2,
\end{equation}
and
\begin{equation}\label{eq:jointpoly}
 (19568+810\theta)QD_*
 +120000N(\theta,Y)\pm300000\mathcal N_\triangle>0,
 \qquad 0\le\theta\le1/3.
\end{equation}

The degree-\(n\) interval Bernstein coefficient of \(\theta^p\) on
\([0,1/b]\), at index \(i\), is
\(b^{-p}\binom ip/\binom np\).
For \(u^q\) on \([0,1]\), use \(\binom jq/\binom mq\).
On the degree-\(m\) simplex, the coefficient of
\(\alpha_1^{q_1}\alpha_2^{q_2}\alpha_3^{q_3}\) at
\(\gamma_1+\gamma_2+\gamma_3=m\) is
\[
 \frac{\prod_{\ell=1}^3
 \gamma_\ell(\gamma_\ell-1)\cdots(\gamma_\ell-q_\ell+1)}
 {m(m-1)\cdots(m-q_1-q_2-q_3+1)}.
\]
Empty ratios equal one and impossible indices give zero.
Linearity and products give the coefficients on product domains.

For the rectangle, let \(A^{(2)}_{ij}\) and \(d_i^{(2)}\) be the
coefficients of \(N(\theta,Y)\) and \(D_*\) at degrees \((38,19)\)
and \(38\), on \(\theta\in[0,1/2]\).
For the joint domain, let \(A^{(3)}_{ij},B_{i\gamma},d_i,t_i\) be
the coefficients of \(N(\theta,Y),\mathcal N_\triangle,D_*,\theta D_*\)
at the degrees in Table~\ref{tab:bernstein}. Put
\[
 L_i=
 Q(19568d_i+810t_i)
 +120000\min_jA^{(3)}_{ij}
 -300000\max_\gamma|B_{i\gamma}|.
\]
The coefficient inequalities in Table~\ref{tab:bernstein}, obtained
by the displayed rational conversion formulas, imply
\eqref{eq:rectangle} and both signs of \eqref{eq:jointpoly}.
For the joint domain the normalized lower margin is greater than
\(1/6000\).

\begin{table}[ht]
\centering
\small
\caption{Bernstein data. The simplex degree refers to
\(\alpha_1+\alpha_2+\alpha_3=1\).}
\label{tab:bernstein}
\begin{tabular}{@{}p{.19\textwidth}p{.32\textwidth}p{.39\textwidth}@{}}
\toprule
Domain & Degrees and indices & Coefficient lower bounds\\
\midrule
Rectangle:
\(0\le\theta\le1/2\) &
\(\deg_\theta=38,\ \deg_u=19\);
\(0\le i\le38,\ 0\le j\le19\) &
\(100A^{(2)}_{ij}+3Qd_i^{(2)}>0\)
(\(780\) coefficients)\\[3pt]
Joint:
\(0\le\theta\le1/3\) &
\(\deg_\theta=39,\ \deg_u=19,\ \deg_{\rm simplex}=19\);
\(0\le i\le39\) &
\(L_i>20Qd_i\)
(\(40\) joint bounds)\\
\bottomrule
\end{tabular}
\end{table}

Dividing \eqref{eq:jointpoly} by \(120000QD_*\) gives
\[
 \frac16-\varepsilon_0(16-30\theta)
 +\mathcal E_\theta^{[10]}(Y)
 \pm\frac52\Omega_\theta^{[10]}.
\]
The superscript denotes the ten-term replacement from
Table~\ref{tab:cosecant}.
By \eqref{eq:uniformtail}, the omitted terms have total magnitude
at most
\(\varepsilon_0+15(1-2\theta)\varepsilon_0\).
This proves \eqref{eq:continuouslemma}.
Likewise, \eqref{eq:rectangle} gives
\begin{equation}\label{eq:rectanglebound}
 \mathcal E_\theta(y)>-\frac3{100}-\varepsilon_0,
 \qquad
 0\le\theta\le1/2,\quad\theta\le y\le1-\theta.
\end{equation}

\begingroup
\small
\bibliographystyle{unsrt}
\bibliography{cglmp-references}
\endgroup

\end{document}